\documentclass[11pt]{article}

% -------------------------------------------------
% Packages
% -------------------------------------------------
\usepackage[T1]{fontenc}
\usepackage[utf8]{inputenc}
\usepackage[english]{babel}
\usepackage[margin=.92in]{geometry}
\usepackage{microtype}
\usepackage[numbers]{natbib}
\usepackage{xr}
\usepackage[colorlinks=true,linkcolor=blue,citecolor=blue,urlcolor=blue]{hyperref}
\usepackage{url}

\usepackage{amsmath,amssymb,amsthm,mathtools,mathrsfs}
\usepackage{array,booktabs,tabularx,makecell,longtable}
\usepackage{graphicx,subcaption,float}
\usepackage{xcolor}
\usepackage{enumitem}
\usepackage{algorithm}
\usepackage{algpseudocode}
\usepackage{cancel}
\usepackage{placeins}

% -------------------------------------------------
% Theorem environments
% -------------------------------------------------
\theoremstyle{plain}
\newtheorem{thm}{Theorem}[section]
\newtheorem{cor}[thm]{Corollary}
\newtheorem{lem}[thm]{Lemma}
\newtheorem{prop}[thm]{Proposition}

\theoremstyle{definition}
\newtheorem{defn}[thm]{Definition}
\newtheorem{assumption}[thm]{Assumption}
\newtheorem{ex}[thm]{Example}

\theoremstyle{remark}
\newtheorem{rem}[thm]{Remark}

% -------------------------------------------------
% Commands
% -------------------------------------------------
\DeclarePairedDelimiter{\nint}{\lfloor}{\rceil}
\DeclarePairedDelimiter{\abs}{\lvert}{\rvert}

\newcommand{\R}{\mathbb{R}}
\newcommand{\EE}{\mathbb{E}}
\newcommand{\PP}{\mathbb{P}}
\newcommand{\op}{\mathrm{op}}
\newcommand{\tr}{\operatorname{tr}}
\newcommand{\rank}{\operatorname{rank}}
\newcommand{\card}{\operatorname{Card}}
\newcommand{\acc}{\operatorname{acc}}
\newcommand{\1}{\mathbf{1}}
\newcommand{\de}{\partial}
\newcommand{\eps}{\varepsilon}
\newcommand{\di}{\mathrm{d}}
\newcommand{\proofloc}[1]{\par\noindent\emph{Proof.} #1\par}

\newcommand{\maybeincludegraphics}[2][]{%
  \IfFileExists{#2}{\includegraphics[#1]{#2}}{%
  \fbox{\parbox[c][2.2in][c]{0.92\linewidth}{\centering Missing figure: \detokenize{#2}}}}}

% -------------------------------------------------
% Title
% -------------------------------------------------
\title{\textbf{Online Resource 1: Proofs and Mathematical Supplement}\\[4pt]
\large Supplementary information for \emph{Pruning Deep Neural Networks via the Marchenko--Pastur Distribution}}

\author{
    Leonid Berlyand \\
    {\small Department of Mathematics} \\
    {\small Pennsylvania State University} \\
    {\small University Park, PA 16802, USA}
    \and
    Theo Bourdais\\
    {\small Department of Computing and Mathematical Sciences} \\
    {\small California Institute of Technology} \\
    {\small Pasadena, CA 91125, USA}
    \and
    Houman Owhadi\\
    {\small Department of Computing and Mathematical Sciences} \\
    {\small California Institute of Technology} \\
    {\small Pasadena, CA 91125, USA}
    \and
    Yitzchak Shmalo\thanks{Corresponding author: Yitzchak Shmalo. Email: \texttt{yitzchak.shmalo@gmail.com}.} \\
    {\small Department of Mathematics} \\
    {\small Pennsylvania State University} \\
    {\small University Park, PA 16802, USA}
}
\date{}

\newcommand{\mainpaper}{the main paper}
\newcommand{\proofresource}{Online Resource 1}
\newcommand{\numericresource}{Online Resource 2}

% Main-manuscript labels are imported for theorem references.  We suppress
% imported bibliography metadata so this supplementary PDF has its own clean
% bibliography pass.
\makeatletter
\let\esmoldbibcite\bibcite
\let\bibcite\@gobbletwo
\externaldocument{main}
\let\bibcite\esmoldbibcite
\makeatother

\begin{document}
\maketitle

\noindent\textbf{Article title:} Pruning Deep Neural Networks via the Marchenko--Pastur Distribution\\
\textbf{Corresponding author:} Yitzchak Shmalo, \texttt{yitzchak.shmalo@gmail.com}

\medskip
\noindent\textbf{Caption.} This file is Online Resource 1. It contains the proof details, Gaussian and abstract random-matrix specializations, finite-rank deformation material, and mathematical corollaries supporting the main manuscript. The main manuscript cites this file as ``Online Resource 1.'' It is intended to be reviewed and published online exactly as submitted.

\tableofcontents

\section{Additional Deterministic Certificate Corollaries}
\label{app:deterministic_relocated_corollaries}

\begin{cor}[Weighted elastic-net mask certificate]
\label{cor:weighted_mask_certificate}
Under the hypotheses of Theorem~\ref{thm:deterministic_mask_certificate}, let \(\Gamma_2=(\gamma_{2,ij})\) and \(\Gamma_1=(\gamma_{1,ij})\) be nonnegative coordinatewise weights, and define
\[
\mathcal J_{\Gamma}(A):=
\mathcal L_{\mathrm{CE}}(A)
+
\sum_{i,j}\gamma_{2,ij}A_{ij}^2
+
\sum_{i,j}\gamma_{1,ij}|A_{ij}|
+
\mathcal C.
\]
Then, for every \(\theta\in[0,1]\),
\[
\mathcal J_{\Gamma}(A_\theta)-\mathcal J_{\Gamma}(A_1)
\le
(1-\theta)B_T(R)
-
(1-\theta^2)\sum_{i,j}\gamma_{2,ij}R_{ij}^2
-
(1-\theta)\sum_{i,j}\gamma_{1,ij}|R_{ij}|.
\]
In particular, full pruning is certified to decrease the weighted elastic-net objective whenever
\[
B_T(R)
<
\sum_{i,j}\gamma_{2,ij}R_{ij}^2
+
\sum_{i,j}\gamma_{1,ij}|R_{ij}|.
\]
\end{cor}

\begin{cor}[Weighted stationarity controls removable masks]
\label{cor:weighted_mask_stationarity}
Under the hypotheses of Corollary~\ref{cor:weighted_mask_certificate}, assume the one-sided directional stationarity condition
\[
\liminf_{h\downarrow0}
\frac{\mathcal J_{\Gamma}(A_{1-h})-\mathcal J_{\Gamma}(A_1)}{h}
\ge
-\varepsilon
\]
holds for some \(\varepsilon\ge0\). Then
\[
2\sum_{i,j}\gamma_{2,ij}R_{ij}^2
+
\sum_{i,j}\gamma_{1,ij}|R_{ij}|
\le
B_T(R)+\varepsilon.
\]
More generally, if
\[
\mathcal J_{\Gamma}(A_\theta)\ge \mathcal J_{\Gamma}(A_1)-\varepsilon(1-\theta)
\qquad\forall\theta\in[0,1],
\]
then for every \(\theta<1\),
\[
(1+\theta)\sum_{i,j}\gamma_{2,ij}R_{ij}^2
+
\sum_{i,j}\gamma_{1,ij}|R_{ij}|
\le
B_T(R)+\varepsilon.
\]
\end{cor}

\begin{cor}[Deterministic multi-layer margin stability]
\label{cor:deterministic_multilayer_margin_stability}
Use the ordered pruning sequence and notation of Theorem~\ref{thm:deterministic_multilayer_mask_certificate}. Let \(z^{(j)}(s)\) be the logit vector after the \(j\)-th mask operation, with \(z^{(0)}\) denoting the original network and \(z^{(M)}\) the final pruned network. Suppose that for each \(s\in T\) and each step \(j\) there is a deterministic bound \(\eta_{j,s}\ge0\) such that
\[
\|z^{(j)}(s)-z^{(j-1)}(s)\|_\infty\le \eta_{j,s}.
\]
Define \(\eta_s^{\mathrm{tot}}:=\sum_{j=1}^M\eta_{j,s}\). Then
\[
\big|\mathcal L_{\mathrm{CE}}(\alpha^{(M)})-\mathcal L_{\mathrm{CE}}(\alpha^{(0)})\big|
\le
\frac{2}{|T|}\sum_{s\in T}\eta_s^{\mathrm{tot}}.
\]
Moreover, if \(p_0(s)\in\arg\max_j z^{(0)}_j(s)\) and
\[
\Delta_0^{\mathrm{pred}}(s):=
z_{p_0(s)}^{(0)}(s)-\max_{j\ne p_0(s)}z_j^{(0)}(s),
\]
then the predicted label at \(s\) can change between the original and final networks only if
\[
\Delta_0^{\mathrm{pred}}(s)\le2\eta_s^{\mathrm{tot}}.
\]
For the true-label margin
\[
\Delta_0^{\mathrm{true}}(s):=
z_{C(s)}^{(0)}(s)-\max_{j\ne C(s)}z_j^{(0)}(s),
\]
the accuracy change obeys
\[
\big|\acc_{\alpha^{(M)}}(T)-\acc_{\alpha^{(0)}}(T)\big|
\le
\frac1{|T|}
\#\{s\in T:|\Delta_0^{\mathrm{true}}(s)|\le2\eta_s^{\mathrm{tot}}\}.
\]
\end{cor}

\begin{cor}[Sub-Gaussian sufficient condition for the data-path budget]
\label{cor:subgaussian_data_path_budget}
Let \(v_s=\psi_1(s)\), and assume \(L_s\) are nonnegative upper bounds for the post-layer Lipschitz constants, deterministic or measurable with respect to variables independent of \(R\). Let the rows \(r_1,\dots,r_n\) of \(R\) be independent and centered conditionally on those variables, and suppose that for some \(c,g>0\),
\[
\PP\!\left(
|\langle r_i,v\rangle|>t\,\middle|\,v
\right)
\le
2\exp\!\left(
-\frac{c n t^2}{g\|v\|_2^2}
\right)
\]
for all \(i\), \(v\), and \(t>0\). Then, with probability at least \(1-\delta\),
\[
B_T(R)
\le
\frac2{|T|}
\sum_{s\in T}
L_s\|v_s\|_2
\sqrt{
\frac{g}{c n}\log\frac{2n|T|}{\delta}
}.
\]
\end{cor}

\begin{cor}[Simultaneous multi-layer Gaussian data-path budget]
\label{cor:multilayer_gaussian_data_path_budget}
Consider the ordered finite pruning sequence in Theorem~\ref{thm:deterministic_multilayer_mask_certificate}. At step \(j\), let the removed component be modeled by \(R_j\in\R^{n_j\times m_j}\) with independent \(N(0,g_j/n_j)\) entries, conditionally independent of the intermediate pre-layer vectors \(v_{j,s}\) and Lipschitz bounds \(L_{j,s}\). Let \(\delta_1,\dots,\delta_M>0\) satisfy \(\sum_{j=1}^M\delta_j\le\delta\). Then, with probability at least \(1-\delta\), all layerwise budgets satisfy
\[
B_j(R_j)
\le
\frac2{|T|}
\sum_{s\in T}
L_{j,s}\|v_{j,s}\|_2
\sqrt{
\frac{2g_j}{n_j}\log\frac{2n_j|T|}{\delta_j}
}
\qquad
(1\le j\le M).
\]
Consequently, on the same event, Theorem~\ref{thm:deterministic_multilayer_mask_certificate} holds with \(\sum_jB_j(R_j)\) replaced by the sum of the displayed right-hand sides.
\end{cor}

\begin{cor}[Simultaneous multi-layer sub-Gaussian data-path budget]
\label{cor:multilayer_subgaussian_data_path_budget}
Use the ordered pruning sequence and notation of Corollary~\ref{cor:multilayer_gaussian_data_path_budget}, but at each step \(j\) replace the Gaussian model by the row-concentration hypothesis of Corollary~\ref{cor:subgaussian_data_path_budget} with \((R,n,c,g,v_s,L_s)\) replaced by \((R_j,n_j,c_j,g_j,v_{j,s},L_{j,s})\). Let \(\delta_1,\dots,\delta_M>0\) satisfy \(\sum_{j=1}^M\delta_j\le\delta\). Then, with probability at least \(1-\delta\),
\[
B_j(R_j)
\le
\frac2{|T|}
\sum_{s\in T}
L_{j,s}\|v_{j,s}\|_2
\sqrt{
\frac{g_j}{c_jn_j}\log\frac{2n_j|T|}{\delta_j}
}
\qquad
(1\le j\le M).
\]
Consequently, the same deterministic multi-layer certificate applies with the displayed sub-Gaussian bounds.
\end{cor}

\section{Notes on the Abstract Additive Framework}
\label{app:abstract_additive_notes}

\subsection{Relocated abstract additive corollaries}
\label{app:abstract_relocated_corollaries}

\begin{cor}[Growing-set abstract loss reduction]
\label{cor:growing_set_abstract_loss}
Let \(T_N\) be a finite labeled set that may depend on \(N\). Let \(L_{\cancel{\mathrm{reg}},T_N}\) denote cross-entropy averaged over \(T_N\), and let \(L_{T_N}\) denote the corresponding regularized objective with the same Frobenius penalty as in \eqref{loss_noise_det_2}. Assume Assumptions~\ref{as4}--\ref{as5},
\[
\langle S,R\rangle_F=o_{\PP}(1),
\qquad
|T_N|d_2(N)\to0,
\]
and
\[
\frac{1}{|T_N|}\sum_{s\in T_N}a_\phi(N,s)\to0
\]
in probability. Then
\[
L_{T_N}(\alpha_S)
=
L_{T_N}(\alpha_W)
-\lambda_{\mathrm{reg}}\|R\|_F^2
+o_{\PP}(1).
\]
\end{cor}

\begin{cor}
\label{cor:scaled_reduction_of_loss_general}
Under the hypotheses of Theorem~\ref{theorem_reduction_of_loss_general}, for any fixed \(\vartheta\in[0,1]\), define \(W(\vartheta)=S+\vartheta R\). Then
\[
L(\alpha_{W(\vartheta)})
=
L(\alpha_W)-\lambda_{\mathrm{reg}}(1-\vartheta^2)\|R\|_F^2+o_{\PP}(1)
\qquad(N\to\infty).
\]
The same asymptotic statement holds along any deterministic sequence \(\vartheta_N\in[0,1]\).
\end{cor}

\begin{cor}[Finite-\(N\) scaled loss bound along the denoising path]
\label{cor:scaled_reduction_of_loss_general_finiteN}
Assume the hypotheses of Theorem~\ref{theorem_reduction_of_loss_general}, and define \(W(\vartheta)=S+\vartheta R\) for \(\vartheta\in[0,1]\). Then for every \(\eta>0\) there exists an event \(E_{N,\eta}\) with
\[
\PP(E_{N,\eta})\ge 1-|T|d_2(N)-\PP\!\big(|\langle S,R\rangle_F|>\eta\big)
\]
such that on \(E_{N,\eta}\),
\[
L(\alpha_{W(\vartheta)})
\le
L(\alpha_W)
-\lambda_{\mathrm{reg}}(1-\vartheta^2)\|R\|_F^2
+(1-\vartheta)\frac{2}{|T|}\sum_{s\in T}a_\phi(N,s)
+2\lambda_{\mathrm{reg}}(1-\vartheta)\eta
\]
for every \(\vartheta\in[0,1]\).
\end{cor}

\begin{cor}[Eventual supercriticality of persistent spike clusters in a BGN-type deformation regime]
\label{cor:general_BGN_persistent_spikes}
Under the hypotheses of Theorem~\ref{main_result_reducing_noise_general}, assume additionally that for each \(N\), \(W^*(N)=S^*(N)+R^*(N)\) lies in a finite-rank deformation regime satisfying all hypotheses of Theorem~\ref{thm:abstract_inverse_D} for the spikes considered below. Assume the regime has deterministic bulk edge \(b_+(N)\), critical threshold \(\theta_c(N)\), strictly decreasing outlier transform
\[
D_N:(b_+(N),\infty)\to (0,D_N(b_+(N)^+)),
\]
local Lipschitz/invertibility intervals around the outliers, and concentration of the relevant sample singular values in those intervals. Assume \(\theta_c(N)\to0\), \(b_+(N)\to0\), and that for some fixed \(k\ge1\) there exist constants \(\bar\sigma_1>\cdots>\bar\sigma_k>0\) such that \(\sigma_i^*(N)\to\bar\sigma_i\) in probability for each \(1\le i\le k\). Then:
\begin{enumerate}[leftmargin=2em]
\item \(\PP(\sigma_i^*(N)>\theta_c(N))\to1\) for each \(1\le i\le k\);
\item these \(k\) spikes are eventually supercritical, so the corresponding singular values of \(W^*(N)\) form outlier clusters above \(b_+(N)\). If the \(i\)-th spike is separated from the other nonzero singular values of \(S^*(N)\) by a gap \(\Delta_i(N)>0\) with \(\|R^*(N)\|_{\op}/\Delta_i(N)\to0\) in probability, then the associated rank-one left and right spectral projectors converge in operator norm to the corresponding projectors of \(S^*(N)\);
\item defining \(\widehat\sigma_i^{(D)}(N):=D_N(s_i(W^*(N)))^{-1/2}\) on the event \(s_i(W^*(N))>b_+(N)\), whose probability tends to one, and arbitrarily off that event, one has \(\widehat\sigma_i^{(D)}(N)-\sigma_i^*(N)\to0\) in probability.
\end{enumerate}
\end{cor}

Assumptions~\ref{as4}--\ref{as5} are the abstract data-path version of the Gaussian conditions in Appendix~\ref{Main_theory}; Assumption~\ref{as2} gives the iid-Gaussian sufficient condition, while the optional ESD condition is used only for spectral conclusions.

\begin{ex}
The square iid-Gaussian case with variance \(1/N\) gives Assumption~\ref{as5} with \(J(v)=\|v\|_2\), \(d_1(N)=2\sqrt{2\log(2N)/N}\), and \(d_2(N)=1/N\); see Corollary~\ref{cor:gaussian_data_path_budget}. The conditional formulation also accommodates sub-Gaussian rows, orthogonally invariant perturbations, or moderate training-induced dependence when the concentration estimate survives.
\end{ex}

\begin{ex}[Sub-Gaussian rows]
\label{ex:subgaussian_rows}
The row-concentration variant in Corollary~\ref{cor:subgaussian_data_path_budget} gives the corresponding Assumption~\ref{as5} sufficient condition.
\end{ex}

Assumption~\ref{as6} is used only for spectral conclusions. Loss reduction uses the data-path smallness \(a_\phi(N,s)\to0\) and the cross-term condition; stationarity-collapse results additionally use the stated one-sided stationarity hypothesis.

\begin{rem}[Regularisation scaling]
\label{rem:lambda_reg_scaling}
Throughout the loss-reduction statements, the elastic-net coefficients \(\lambda_1,\lambda_2\) are treated as specified constants. An iid-Gaussian admissible component \(R\in\mathbb R^{N\times N}\) with variance \(g/N\) has \(\|R\|_F^2=\Theta(gN)\) in expectation, so the loss-reduction term \(\lambda_2\|R\|_F^2\) scales as \(\Theta(\lambda_2 gN)\). The certificate inequality of Theorem~\ref{thm:deterministic_mask_certificate} holds for any specified coefficient scaling, including coefficients normalized by parameter count or layer width. The data-path budget \(B_T(R)\) is independent of \(\lambda_1,\lambda_2\), while the elastic-net reward is not; only the numerical certificate margin changes with regularizer scale.
\end{rem}

\begin{rem}[Beyond pure Frobenius regularization]
\label{rem:bregman_regularizer}
The quadratic penalty in \eqref{loss_noise_det_2} supplies a coercive lower bound of order \(\|R\|_F^2\). For a differentiable \(\mu\)-strongly convex regularizer \(\Omega\),
\[
\Omega(S+R)-\Omega(S)=\langle\nabla\Omega(S),R\rangle+D_\Omega(S+R,S),
\qquad
D_\Omega(S+R,S)\ge\frac{\mu}{2}\|R\|_F^2.
\]
The extension is valid when the first-order term is controlled or negligible. In the Frobenius case this term is \(2\langle S,R\rangle_F\). For elastic net, choose \(\xi\in\partial\|S\|_1\) and assume \(\langle\mu_1\xi+2\mu_2S,R\rangle=o_{\PP}(1)\); the quadratic part supplies the strong convexity.
\end{rem}

The stationarity theorem assumes only one-sided approximate stationarity at \(\vartheta=1\); exact local minimality is sufficient but stronger. Without the no-atom condition \(\PP(\|R^*(N)\|_F=\epsilon)=0\), the same proof gives limsup/liminf threshold bounds rather than exact convergence. Corollary~\ref{cor:general_BGN_persistent_spikes} is the corresponding abstract BBP persistence statement.

Dependency map: deterministic mask certificates use the support condition
\[
S=M\odot A,\qquad R=(1-M)\odot A,
\]
so \(S\odot R=0\). Abstract additive loss statements replace that support condition by the relevant cross-term control, such as \(\langle S,R\rangle_F=o_{\PP}(1)\). Random budget corollaries provide sufficient conditions for \(B_T(R)\).

\begin{rem}[Local Lipschitz constants $L_s$]
\label{rem:lipschitz_local}
The constants \(L_s\) in \eqref{eq:data_path_budget} are local post-layer \(\ell_\infty\)-Lipschitz bounds; for ViTs this paper treats them as architecture-dependent and uses the empirical proxy audit in Online Resource 2.
\end{rem}

\section{Gaussian Specialization and Spectral Interpretation}
\label{Main_theory}
% =================================================

The deterministic mask certificates in Section~\ref{Main_theory_gen} are closest to the implemented unstructured pruning operation; this appendix gives Gaussian sufficient conditions and spectral interpretations for admissible random-like components.

\subsection{Assumptions for the Gaussian specialization}
\label{assumptions}

The finite-\(T\) convention from Section~\ref{sec:abstract_additive_extensions} remains in force unless stated otherwise.

Whenever \(\lambda_+\) denotes the right MP edge for the eigenvalues of
\[
X_\ell=\frac1N W_\ell^\top W_\ell,
\]
we write
\[
\sigma_+ := \sqrt{N\lambda_+}
\]
for the corresponding edge on the singular-value scale.

\begin{assumption}[Architecture and outer-layer scaling]
\label{as1}
We consider the simplified three-layer network
\begin{align}
X(\alpha,s) &= \lambda\circ W_3 \circ \lambda\circ W_2 \circ \lambda\circ W_1 s,\qquad s\in T,\\
\phi(\alpha,s) &= \rho\circ X(\alpha,s),
\end{align}
where \(\lambda\) is the coordinatewise absolute value or ReLU, and \(W_1,W_2,W_3\) are the weight matrices. Biases are omitted for simplicity.

For fixed training time \(t\), define
\begin{equation}
\label{eqans}
a(N,s):=\|W_3\|_{\infty}\,\|W_1s\|_2\sqrt{\frac{\log(2N)}{N}}.
\end{equation}
We assume that for every \(s\in T\), \(a(N,s)\to 0\) as \(N\to\infty\). The same assumption is imposed at any admissible limit point considered below.
\end{assumption}

\begin{rem}[Interpretation of Assumption~\ref{as1}]
This is the concrete version of the abstract factorization Assumption~\ref{as4}.  The target layer is \(W_2\); the computations before \(W_2\) play the role of \(\psi_1\), and the computations after \(W_2\) play the role of \(\psi_2\).  The quantity \(a(N,s)\) is the Gaussian-model version of \(a_\phi(N,s)=h_\phi(s)d_1(N)\): it is the logit perturbation scale predicted by multiplying an \(N^{-1/2}\sqrt{\log N}\)-size random effect through the outer layer.  The assumption \(a(N,s)\to0\) is therefore the concrete data-path smallness condition needed for the loss and margin bounds.
\end{rem}

\begin{assumption}[Admissible Gaussian decomposition of the middle layer]
\label{as2}
For each training time \(t\), the middle layer \(W_2(t)\in\R^{N\times N}\) admits a decomposition
\[
W_2(t)=R_2(t)+S_2(t),
\]
where:
\begin{enumerate}[leftmargin=2em]
\item \(R_2(t)\) has i.i.d.\ centered Gaussian entries
\[
R_2(t)_{ij}\overset{\mathrm{i.i.d.}}{\sim}\mathcal N\!\left(0,\frac{g(t)}{N}\right),
\qquad 0\le g(t)\le 1;
\]
where \(g(t)\) is a deterministic variance scale;
\item \(R_2(t)\) is independent of \(S_2(t)\);
\item \(R_2(t)\) is also independent of the outer-layer parameters \(W_1(t)\) and \(W_3(t)\).
\end{enumerate}
\end{assumption}

\begin{rem}
The variance scaling \(g(t)/N\) ensures that \(R_2(t)\) has bounded operator norm as \(N\to\infty\). Rectangular variants \(W_2(t)\in\R^{N\times M}\) with \(M/N\to c\in(0,\infty)\) can be treated similarly.
\end{rem}

\begin{rem}[Interpretation of Assumption~\ref{as2}]
This is a stylized additive-noise model, not a claim that the implemented ViT mask literally samples an independent Gaussian matrix.  It is used to prove a sufficient condition: if the removed component behaves like an independent variance-\(g/N\) random matrix along the data path, then it satisfies the perturbation bounds required by the abstract theory.  In the notation of Assumption~\ref{as5}, it gives \(J(v)=\|v\|_2\), \(d_1(N)\asymp\sqrt{\log N/N}\), and \(d_2(N)\asymp1/N\).  The same variance scale has MP singular-value edge \(2\sqrt g\) in the square case, which is why the fitted MP edge is a natural diagnostic in the algorithms.
\end{rem}

\begin{assumption}[Sub-root-\(N\) structured component for the loss theorems]
\label{as3}
For each training time \(t\), the structured component \(S_2(t)\) satisfies
\[
\|S_2(t)\|_F\le B_N,
\qquad
B_N=o(\sqrt N),
\]
for some positive deterministic sequence \(B_N\) independent of \(t\). Equivalently, the normalized size condition \(\|S_2(t)\|_F^2/N\to0\) is enough. The loss-reduction theorems use this assumption only to prove
\[
\langle S_2(t),R_2(t)\rangle_F=o_{\PP}(1);
\]
one may replace it directly by this cross-term condition whenever it is known by another argument.
\end{assumption}

\begin{rem}[Interpretation of Assumption~\ref{as3}]
This assumption is a convenient Gaussian sufficient condition for the cross-term bound \(\langle S_2,R_2\rangle_F=o_{\PP}(1)\).  Conditional on \(S_2\),
\[
\mathrm{Var}\!\left(\langle S_2,R_2\rangle_F\,\middle|\,S_2\right)
=
\frac{g(t)}{N}\|S_2\|_F^2.
\]
Thus \(B_N=o(\sqrt N)\) forces this variance to vanish; this is the Gaussian counterpart of the explicit cross-term hypothesis used in Section~\ref{sec:abstract_additive_extensions}.
\end{rem}

\begin{assumption}[Approximate low rank for the spectral corollaries]
\label{as3spectral}
Whenever a spectral conclusion is invoked, we additionally assume that the limit structured component \(S_2^*(N)\) is approximately low rank in the sense that there exists a sequence \(k_N=o(N)\) such that
\[
\sum_{i>k_N}s_i\!\big(S_2^*(N)\big)^2\to0
\qquad\text{in probability as }N\to\infty.
\]
The exact fixed-rank regime is the special case \(k_N\equiv r\).
\end{assumption}

\begin{rem}[Interpretation of Assumption~\ref{as3spectral}]
This is the Gaussian specialization of Assumption~\ref{as6} and is used only by the spectral corollaries.
\end{rem}

\begin{rem}[No spectral-separation assumption]
Assumptions~\ref{as3}--\ref{as3spectral} do \emph{not} require the informative singular values of \(S_2\) to lie above the MP edge \(\sigma_+\). Thus signal directions may still be embedded in the bulk created by \(R_2\).
\end{rem}

\begin{defn}[Admissible decomposition]
\label{def:admissible_decomposition}
A decomposition \(W_2=S_2+R_2\) is called \emph{admissible} if it satisfies Assumptions~\ref{as2}--\ref{as3}. At a limit point \(W_2^*(N)\), an admissible decomposition \(W_2^*(N)=S_2^*(N)+R_2^*(N)\) additionally requires that the outer-layer scaling condition from Assumption~\ref{as1} holds for the corresponding limit outer layers \(W_1^*(N),W_3^*(N)\), i.e.
\[
a_*(N,s):=\|W_3^*(N)\|_\infty\|W_1^*(N)s\|_2\sqrt{\frac{\log(2N)}{N}}\to0
\qquad \text{for each } s\in T.
\]
A \emph{spectrally admissible} decomposition is an admissible decomposition that also satisfies Assumption~\ref{as3spectral}.
\end{defn}

\begin{defn}[Local-path convergence of an admissible decomposition]
\label{def:local_path_convergence}
An admissible path decomposition
\[
W_2(t,N)=S_2(t,N)+R_2(t,N)
\]
is said to converge locally to an admissible limit decomposition
\[
W_2^*(N)=S_2^*(N)+R_2^*(N)
\]
if, for each fixed \(N\),
\[
\|S_2(t,N)-S_2^*(N)\|_F+\|R_2(t,N)-R_2^*(N)\|_F\to0
\qquad\text{in probability as }t\to\infty.
\]
\end{defn}

\subsection{Key perturbation lemma}
\label{subsec:key_perturbation_lemma}

\begin{lem}
\label{lem:remove_R}
Suppose Assumptions~\ref{as1} and \ref{as2} hold. Define
\[
\tau_N(s):=
2\sqrt{2}\,\|W_3\|_\infty \|W_1s\|_2\sqrt{\frac{\log(2N)}{N}}
=2\sqrt{2}\,a(N,s).
\]
If we replace \(W_2=S_2+R_2\) by \(S_2\), then
\[
\PP\!\left(
\left\|X(s,\alpha_{S_2})-X(s,\alpha_{W_2})\right\|_\infty
\le \tau_N(s)
\right)
\ge 1-\frac{1}{N}.
\]
\end{lem}
\proofloc{See Appendix~\ref{app:proof_perturb}, Proof of Lemma~\ref{lem:remove_R}.}

For later use, define the classification margin on a finite test set \(T'\):
\[
\delta X(s,\alpha):=
X_{C(s)}(s,\alpha)-\max_{j\neq C(s)}X_j(s,\alpha),
\]
and the corresponding accuracy
\[
\acc_\alpha(T'):=\frac1{|T'|}\#\{s\in T':\delta X(s,\alpha)>0\}.
\]

\begin{cor}[Margin-sensitive label stability]
\label{cor:margin_accuracy}
Suppose Assumptions~\ref{as1} and \ref{as2} hold. Let \(T'\) be a finite test set. Then
\[
\PP\!\left(
\big|\acc_{\alpha_{W_2}}(T')-\acc_{\alpha_{S_2}}(T')\big|
\le
\frac1{|T'|}
\#\Big\{
s\in T': |\delta X(s,\alpha_{W_2})|\le 2\tau_N(s)
\Big\}
\right)
\ge
1-\frac{|T'|}{N}.
\]
In particular, if
\[
\min_{s\in T'} |\delta X(s,\alpha_{W_2})|
>
2\max_{s\in T'}\tau_N(s),
\]
then
\[
\acc_{\alpha_{W_2}}(T')=\acc_{\alpha_{S_2}}(T')
\]
with probability at least \(1-\frac{|T'|}{N}\).
\end{cor}
\proofloc{See Appendix~\ref{app:proofs_main}, Proof of Corollary~\ref{cor:margin_accuracy}.}

\subsection{A theorem on how removing randomness reduces the regularized loss}

We work with the \(\ell_2\)-regularized loss
\begin{equation}
\label{loss_noise_det_1}
L(\alpha(t))
=
-\frac{1}{|T|}\sum_{s\in T}\log\big(\phi_{C(s)}(s,\alpha(t))\big)
+
\lambda_{\mathrm{reg}}
\sum_{i=1}^{L}\|W_i(t)\|_F^2,
\qquad \lambda_{\mathrm{reg}}>0.
\end{equation}

We first isolate the unregularized perturbation term.

\begin{lem}
\label{lem:remove_R_loosandacc}
Suppose Assumptions~\ref{as1}--\ref{as2} hold. Define the unregularized cross-entropy
\[
L_{\cancel{\mathrm{reg}}}(\alpha)
:=
-\frac{1}{|T|}\sum_{s\in T}\log\big(\phi_{C(s)}(s,\alpha)\big).
\]
Then there exists \(N_0\) such that for all \(N\ge N_0\),
\[
\PP\!\left(
\big|L_{\cancel{\mathrm{reg}}}(\alpha_{S_2})-L_{\cancel{\mathrm{reg}}}(\alpha_{W_2})\big|
\le
\frac{4\sqrt{2}}{|T|}\sum_{s\in T} a(N,s)
\right)
\ge
1-\frac{|T|}{N}.
\]
In particular,
\[
L_{\cancel{\mathrm{reg}}}(\alpha_{S_2})-L_{\cancel{\mathrm{reg}}}(\alpha_{W_2})
\to 0
\qquad\text{in probability as }N\to\infty.
\]
\end{lem}
\proofloc{See Appendix~\ref{app:proof_perturb}, Proof of Lemma~\ref{lem:remove_R_loosandacc}.}

\begin{thm}[Confidence-parameter finite-sample perturbation bound]
\label{thm:finite_confidence_gaussian}
Suppose Assumptions~\ref{as1} and \ref{as2} hold, and fix a finite set \(T_0\) and \(\delta\in(0,1)\). Define
\[
\tau_{N,\delta}(s)
:=
\|W_3\|_\infty\|W_1s\|_2
\sqrt{\frac{2g(t)\log(2N|T_0|/\delta)}{N}}.
\]
Then, with probability at least \(1-\delta\),
\[
\left\|X(s,\alpha_{S_2})-X(s,\alpha_{W_2})\right\|_\infty
\le \tau_{N,\delta}(s)
\qquad\forall s\in T_0.
\]
Consequently, for \(T_0=T\),
\[
\left|
L_{\cancel{\mathrm{reg}}}(\alpha_{S_2})
-
L_{\cancel{\mathrm{reg}}}(\alpha_{W_2})
\right|
\le
\frac{2}{|T|}\sum_{s\in T}\tau_{N,\delta}(s)
\]
with probability at least \(1-\delta\).
If Assumption~\ref{as3} also holds, then for every \(\eta>0\),
\[
\PP\!\left(
L(\alpha_{S_2})
\le
L(\alpha_{W_2})
-\lambda_{\mathrm{reg}}\|R_2\|_F^2
+\frac{2}{|T|}\sum_{s\in T}\tau_{N,\delta}(s)
+2\lambda_{\mathrm{reg}}\eta
\right)
\ge
1-\delta-2e^{-N\eta^2/(2B_N^2)}.
\]
\end{thm}
\proofloc{See Appendix~\ref{app:proof_perturb}, Proof of Theorem~\ref{thm:finite_confidence_gaussian}.}

\begin{cor}[Growing-set Gaussian loss bound]
\label{cor:growing_set_gaussian_loss}
Suppose Assumptions~\ref{as1}--\ref{as2} hold with a finite labeled set \(T_N\) that may depend on \(N\), and let \(\delta_N\in(0,1)\) with \(\delta_N\to0\). Let \(L_{\cancel{\mathrm{reg}},T_N}\) denote cross-entropy averaged over \(T_N\), and let \(L_{T_N}\) denote the corresponding regularized objective with the same Frobenius penalty as in \eqref{loss_noise_det_1}. Define
\[
b_{N,\delta_N}(T_N):=
\frac{2}{|T_N|}
\sum_{s\in T_N}
\|W_3\|_\infty\|W_1s\|_2
\sqrt{\frac{2g(t)\log(2N|T_N|/\delta_N)}{N}}.
\]
If \(b_{N,\delta_N}(T_N)\to0\) in probability, then the cross-entropy loss averaged over \(T_N\) satisfies
\[
L_{\cancel{\mathrm{reg}},T_N}(\alpha_{S_2})
-
L_{\cancel{\mathrm{reg}},T_N}(\alpha_{W_2})
\to0
\qquad\text{in probability.}
\]
If Assumption~\ref{as3} also holds, then
\[
L_{T_N}(\alpha_{S_2})
\le
L_{T_N}(\alpha_{W_2})
-\lambda_{\mathrm{reg}}\|R_2\|_F^2
+o_{\PP}(1).
\]
\end{cor}
\proofloc{The bound follows by applying Theorem~\ref{thm:finite_confidence_gaussian} to the finite set \(T_N\) with failure probability \(\delta_N\); the stated assumptions \(\delta_N\to0\) and \(b_{N,\delta_N}(T_N)\to0\) give the result.}

\begin{thm}
\label{theorem_reduction_of_loss}
Suppose Assumptions~\ref{as1}--\ref{as3} hold. If \(W_2=R_2+S_2\) and we replace \(W_2\) by \(S_2\), then for every fixed training time \(t\) and every \(\eta>0\),
\[
\PP\!\left(
L(\alpha_{S_2})
\le
L(\alpha_{W_2})
-\lambda_{\mathrm{reg}}\|R_2\|_F^2
+\frac{4\sqrt{2}}{|T|}\sum_{s\in T}a(N,s)
+2\lambda_{\mathrm{reg}}\eta
\right)
\ge
1-\frac{|T|}{N}-2e^{-N\eta^2/(2B_N^2)}.
\]
In particular,
\[
L(\alpha_{S_2})
=
L(\alpha_{W_2})
-\lambda_{\mathrm{reg}}\|R_2\|_F^2
+o_{\PP}(1)
\qquad(N\to\infty).
\]
\end{thm}
\proofloc{See Appendix~\ref{app:proofs_main}, Proof of Theorem~\ref{theorem_reduction_of_loss}.}

\begin{cor}
\label{cor:scaled_reduction_of_loss}
Suppose Assumptions~\ref{as1}--\ref{as3} hold. For any \(\vartheta\in[0,1]\), define
\[
W_2(\vartheta)=S_2+\vartheta R_2.
\]
Then for every fixed training time \(t\) and every \(\eta>0\), there is an event \(E_{N,\eta}\) with
\[
\PP(E_{N,\eta})
\ge
1-\frac{|T|}{N}-2e^{-N\eta^2/(2B_N^2)}
\]
such that on \(E_{N,\eta}\), for every \(\vartheta\in[0,1]\),
\[
L(\alpha_{W_2(\vartheta)})
\le
L(\alpha_{W_2})
-\lambda_{\mathrm{reg}}(1-\vartheta^2)\|R_2\|_F^2
+\frac{4\sqrt{2}(1-\vartheta)}{|T|}\sum_{s\in T}a(N,s)
+2\lambda_{\mathrm{reg}}(1-\vartheta)\eta.
\]
In particular, for each fixed \(\vartheta\in[0,1]\),
\[
L(\alpha_{W_2(\vartheta)})
=
L(\alpha_{W_2})
-\lambda_{\mathrm{reg}}(1-\vartheta^2)\|R_2\|_F^2
+o_{\PP}(1)
\qquad(N\to\infty).
\]
The same asymptotic statement holds along any deterministic sequence \(\vartheta_N\in[0,1]\).
\end{cor}
\proofloc{See Appendix~\ref{app:proofs_main}, Proof of Corollary~\ref{cor:scaled_reduction_of_loss}.}

\begin{thm}[Entrywise thresholding in a sparse signal-plus-noise model]
\label{thm:entrywise_threshold_support}
Let \(W_N=S_N+R_N\in\R^{N\times M_N}\), where \(R_{ij}\) are independent \(N(0,g_N/N)\) random variables, \(g_N\ge0\), and \(S_N\) is deterministic or independent of \(R_N\). Fix \(\delta\in(0,1)\) and set
\[
t_{N,\delta}
:=
\sqrt{\frac{2g_N\log(2NM_N/\delta)}{N}}.
\]
Define the entrywise hard-thresholded matrix
\[
\mathcal H_t(W)_{ij}:=W_{ij}\,\mathbf 1_{\{|W_{ij}|>t\}}.
\]
Then, with probability at least \(1-\delta\),
\[
\{(i,j): |(S_N)_{ij}|>2t_{N,\delta}\}
\subseteq
\operatorname{supp}\big(\mathcal H_{t_{N,\delta}}(W_N)\big)
\subseteq
\operatorname{supp}(S_N).
\]
Equivalently, thresholding creates no false positives and can miss only entries of \(S_N\) whose magnitude is at most \(2t_{N,\delta}\):
\[
\operatorname{supp}(S_N)\setminus
\operatorname{supp}\big(\mathcal H_{t_{N,\delta}}(W_N)\big)
\subseteq
\{(i,j):0< |(S_N)_{ij}|\le2t_{N,\delta}\}.
\]
Let \(\mathcal B_{N,\delta}\) be the beta-min event
\[
\mathcal B_{N,\delta}:=
\left\{
\min_{(i,j)\in\operatorname{supp}(S_N)} |(S_N)_{ij}|>2t_{N,\delta}
\right\},
\]
with the condition interpreted as vacuous if \(\operatorname{supp}(S_N)=\emptyset\). If \(\PP(\mathcal B_{N,\delta})\ge1-\gamma\), then entrywise thresholding \(W_N\) at level \(t_{N,\delta}\) recovers the support of \(S_N\) exactly with probability at least \(1-\delta-\gamma\):
\[
\operatorname{supp}\big(\mathcal H_{t_{N,\delta}}(W_N)\big)
=
\operatorname{supp}(S_N).
\]
In particular, if \(\mathcal B_{N,\delta}\) holds deterministically, or almost surely when \(S_N\) is random, the recovery probability is at least \(1-\delta\).
This result is a stylized support-recovery statement. Without beta-min separation it still gives a no-false-positive and large-signal-retention certificate, but exact support recovery requires all nonzero structured entries to clear the \(2t_{N,\delta}\) noise buffer. It is not a proof that the ViT pruning masks identify the true structured support.
\end{thm}
\proofloc{See Appendix~\ref{app:proof_perturb}, Proof of Theorem~\ref{thm:entrywise_threshold_support}.}

\subsection{Asymptotic magnitude of admissible noise}

We now replace exact local minimality by a weaker one-sided stationarity condition along the denoising direction \(R_2^*(N)\).

\begin{thm}
\label{main_result_reducing_noise}
Suppose Assumptions~\ref{as1}--\ref{as3} hold along training. Assume that:
\begin{enumerate}[leftmargin=2em]
\item for each \(N\), the parameter vector \(\alpha(t,N)\) converges in probability as \(t\to\infty\) to a limit \(\alpha^*(N)\);
\item the limiting middle layer \(W_2^*(N)\) admits an admissible decomposition
\[
W_2^*(N)=S_2^*(N)+R_2^*(N)
\]
in the sense of Definition~\ref{def:admissible_decomposition}.
\end{enumerate}
\noindent For \(\vartheta\in[0,1]\), let \(\alpha_\vartheta^*(N)\) denote the parameter vector obtained from \(\alpha^*(N)\) by replacing \(W_2^*(N)\) with \(S_2^*(N)+\vartheta R_2^*(N)\). Assume further that there exists a sequence of nonnegative random variables \(\epsilon_N=o_{\PP}(1)\) such that the one-sided directional stationarity event
\[
\liminf_{h\downarrow0}
\frac{
L(\alpha_{1-h}^*(N))-L(\alpha^*(N))
}{h}
\ge -\epsilon_N
\]
has probability tending to one.
Define
\[
\bar a_N:=\frac{4\sqrt{2}}{|T|}\sum_{s\in T}a_*(N,s).
\]
Then
\[
\|R_2^*(N)\|_F^2
\le
\frac{\bar a_N+\epsilon_N}{2\lambda_{\mathrm{reg}}}
+o_{\PP}(1).
\]
In particular, if \(\bar a_N+\epsilon_N\to0\) in probability, then for every \(\epsilon>0\),
\[
\lim_{N\to\infty}\PP\big(\|R_2^*(N)\|_F\le \epsilon\big)=1.
\]

If, in addition, there exists an admissible path decomposition
\[
W_2(t,N)=S_2(t,N)+R_2(t,N)
\]
which converges locally to the limit decomposition in the sense of Definition~\ref{def:local_path_convergence}, then for every \(\epsilon>0\),
\[
\lim_{N\to\infty}\lim_{t\to\infty}\PP\big(\|R_2(t,N)\|_F\le \epsilon\big)=1.
\]
\end{thm}
\proofloc{See Appendix~\ref{app:proofs_main}, Proof of Theorem~\ref{main_result_reducing_noise}.}

\begin{rem}
Theorem~\ref{main_result_reducing_noise} requires only one-sided directional stationarity at \(\vartheta=1\); a genuine local minimum gives \(\epsilon_N\equiv0\). Its second statement transfers the limit-point conclusion to any admissible training-path decomposition converging locally to the limit. Since the Gaussian \(\|R_2^*(N)\|_F\) has no atom at any \(\epsilon>0\), pathwise convergence in probability yields convergence of the corresponding distribution functions at those thresholds.
\end{rem}

\begin{cor}[Variance-scale collapse at the limit point]
\label{cor:variance_scale_collapse}
Under the hypotheses of Theorem~\ref{main_result_reducing_noise}, write the admissible Gaussian limit perturbation as
\[
R_2^*(N)_{ij}\overset{\mathrm{i.i.d.}}{\sim}
\mathcal N\!\left(0,\frac{g_*(N)}{N}\right).
\]
Here \(g_*(N)\) is the deterministic variance scale of the limit perturbation.
Then
\[
g_*(N)\,N\to 0.
\]
In particular, in the square case the MP singular-value edge of the admissible limit perturbation satisfies
\[
\sigma_{+,*}^{\mathrm{bulk}}(N)=2\sqrt{g_*(N)}\to 0.
\]
For a rectangular \(N\times M_N\) perturbation with entries \(N(0,g_*(N)/N)\), the corresponding Frobenius collapse condition is \(g_*(N)M_N\to0\); when \(M_N/N\to c\in(0,\infty)\), this is equivalent to \(g_*(N)N\to0\).
\end{cor}
\proofloc{See Appendix~\ref{app:proofs_main}, Proof of Corollary~\ref{cor:variance_scale_collapse}.}

\begin{cor}[Bulk collapse and spike stabilization at the limit point]
\label{cor:bulk_collapse_spike_stabilization}
Under the hypotheses of Theorem~\ref{main_result_reducing_noise}, assume additionally that the limit decomposition is spectrally admissible in the sense of Definition~\ref{def:admissible_decomposition}. Let
\[
s_1\big(W_2^*(N)\big)\ge \cdots \ge s_N\big(W_2^*(N)\big)
\]
denote the singular values of \(W_2^*(N)\), and let
\[
\sigma_1^*(N)\ge \sigma_2^*(N)\ge \cdots \ge 0
\]
denote the singular values of \(S_2^*(N)\). Let \(k_N=o(N)\) be as in Assumption~\ref{as3spectral}.
Then:
\begin{enumerate}[leftmargin=2em]
\item \(\|R_2^*(N)\|_{\op}\to 0\) in probability;
\item for each fixed \(i\ge1\),
\[
\big|s_i(W_2^*(N))-\sigma_i^*(N)\big|\to 0
\qquad\text{in probability};
\]
\item
\[
\sum_{i>k_N} s_i(W_2^*(N))^2\to 0
\]
in probability;
\item for every \(\epsilon>0\),
\[
\nu_{W_2^*(N)}((\epsilon,\infty))\to 0
\qquad\text{in probability},
\]
where \(\nu_{W_2^*(N)}\) is the singular-value empirical spectral distribution.
\end{enumerate}
If, in addition, \(\sigma_i^*(N)\to \bar\sigma_i\) for some constants \(\bar\sigma_i\ge 0\), then
\[
s_i(W_2^*(N))\to \bar\sigma_i
\qquad\text{in probability for each fixed }i\ge1.
\]
\end{cor}
\proofloc{See Appendix~\ref{app:proofs_main}, Proof of Corollary~\ref{cor:bulk_collapse_spike_stabilization}.}

\begin{cor}[Singular-subspace stabilization under an asymptotic gap condition]
\label{cor:subspace_stabilization}
Under the hypotheses of Theorem~\ref{main_result_reducing_noise}, let
\[
\sigma_1^*(N)\ge \sigma_2^*(N)\ge \cdots \ge 0
\]
denote the singular values of \(S_2^*(N)\). Fix \(k\ge1\), and define
\[
\delta_k(N):=
\sigma_k^*(N)-\sigma_{k+1}^*(N),
\]
with the convention \(\sigma_j^*(N)=0\) for \(j\) beyond the rank of \(S_2^*(N)\). Assume that \(\PP(\delta_k(N)>0)\to1\) and that, after setting the ratio below to \(+\infty\) on the event \(\{\delta_k(N)=0\}\),
\[
\frac{\|R_2^*(N)\|_{\op}}{\delta_k(N)}\to 0
\qquad\text{in probability as }N\to\infty.
\]
Let \(U_{*,k}(N)\) and \(V_{*,k}(N)\) be orthonormal bases for the top-\(k\) left and right singular subspaces of \(S_2^*(N)\), respectively.
Let \(\widehat U_k(N),\widehat V_k(N)\in\R^{N\times k}\) be matrices whose columns are orthonormal left and right singular vectors corresponding to the top \(k\) singular values of \(W_2^*(N)\), and define the orthogonal projectors
\[
P_{U_*,k}(N):=U_{*,k}(N)U_{*,k}(N)^\top,
\qquad
P_{V_*,k}(N):=V_{*,k}(N)V_{*,k}(N)^\top,
\]
\[
\widehat P_{U,k}(N):=\widehat U_k(N)\widehat U_k(N)^\top,
\qquad
\widehat P_{V,k}(N):=\widehat V_k(N)\widehat V_k(N)^\top.
\]
Then
\[
\|\widehat P_{U,k}(N)-P_{U_*,k}(N)\|_{\op}\to0,
\qquad
\|\widehat P_{V,k}(N)-P_{V_*,k}(N)\|_{\op}\to0
\]
in probability as \(N\to\infty\).
\end{cor}
\proofloc{See Appendix~\ref{app:proofs_main}, Proof of Corollary~\ref{cor:subspace_stabilization}.}

\begin{rem}
Corollary~\ref{cor:subspace_stabilization} converts spectral shrinkage into geometric stability: if residual operator-norm noise is asymptotically smaller than \(\delta_k(N)\), the corresponding left/right singular subspaces become identifiable. This does not require a fixed spike floor; \(\delta_k(N)\) may tend to \(0\) if it still dominates the noise.
\end{rem}

\begin{cor}[Eventual Gaussian supercriticality of persistent spikes]
\label{cor:gaussian_BBP_persistent_spikes}
Under the hypotheses of Theorem~\ref{main_result_reducing_noise}, write
\[
R_2^*(N)_{ij}\overset{\mathrm{i.i.d.}}{\sim}
\mathcal N\!\left(0,\frac{g_*(N)}{N}\right).
\]
Assume additionally that, for this corollary only, the admissible limit decomposition lies in the square Gaussian finite-rank deformation regime of Appendix~\ref{app:proofs_main}; in particular, \(S_2^*(N)\) has rank at most \(r_0\) for some deterministic \(r_0\) independent of \(N\).
Assume that for some fixed \(k\ge1\) there exist constants
\[
\bar\sigma_1>\cdots>\bar\sigma_k>0
\]
such that
\[
\sigma_i^*(N)\to \bar\sigma_i
\qquad\text{in probability as }N\to\infty
\]
for each \(1\le i\le k\). Then:
\begin{enumerate}[leftmargin=2em]
\item
\[
\PP\!\big(\sigma_i^*(N)>\sqrt{g_*(N)}\big)\to1
\qquad\text{for each }1\le i\le k;
\]
\item in the sense of Theorem~\ref{thm:BGN_square_variable}, these \(k\) spikes are eventually supercritical, so the corresponding singular values of \(W_2^*(N)\) form BBP outlier clusters above the Gaussian bulk edge \(2\sqrt{g_*(N)}\). Moreover, if the \(i\)-th spike is separated from the other nonzero singular values of \(S_2^*(N)\) by a gap \(\Delta_i(N)>0\) with
\[
\frac{\|R_2^*(N)\|_{\op}}{\Delta_i(N)}\to0
\qquad\text{in probability},
\]
then the rank-one left and right spectral projectors of \(W_2^*(N)\) associated with this spike converge in operator norm to the corresponding projectors of \(S_2^*(N)\);
\item defining the inverse-BBP estimator on the event \(s_i(W_2^*(N))^2>4g_*(N)\), whose probability tends to one, by
\[
\widehat\sigma_i^{\mathrm{BBP}}(N)
:=
\frac12\left(
s_i(W_2^*(N))
+
\sqrt{s_i(W_2^*(N))^2-4g_*(N)}
\right),
\]
and defining it arbitrarily off that event,
one has
\[
\widehat\sigma_i^{\mathrm{BBP}}(N)-\sigma_i^*(N)\to0
\qquad\text{in probability}
\]
for each \(1\le i\le k\).
\end{enumerate}
\end{cor}
\proofloc{See Appendix~\ref{app:proofs_main}, Proof of Corollary~\ref{cor:gaussian_BBP_persistent_spikes}.}

\begin{rem}
Corollary~\ref{cor:gaussian_BBP_persistent_spikes} adds persistent-spike information to bulk collapse: a macroscopic \(\bar\sigma_i>0\) survives as \(g_*(N)\to0\), eventually exceeds the BBP threshold, and, under the stated gap condition, keeps its spectral direction; the inverse-BBP map then estimates the spike strength consistently.
\end{rem}

\begin{rem}[Assumption map for the Gaussian section]
Among the theorem-like results proved in the Gaussian section, the dependence on the hypotheses is as follows:
\begin{enumerate}[leftmargin=2em]
\item Theorem~\ref{RMT_MP_theorem} and Proposition~\ref{prop:mp_bulk_interpretation} are standalone random-matrix statements and do not rely on Assumptions~\ref{as1}--\ref{as3}.
\item Lemma~\ref{lem:remove_R} and Corollary~\ref{cor:margin_accuracy} use only Assumptions~\ref{as1}--\ref{as2}.
\item The first part of Theorem~\ref{thm:finite_confidence_gaussian} uses only Assumptions~\ref{as1}--\ref{as2}.
\item Theorem~\ref{theorem_reduction_of_loss}, Corollary~\ref{cor:scaled_reduction_of_loss}, and the regularized part of Theorem~\ref{thm:finite_confidence_gaussian} use Assumptions~\ref{as1}--\ref{as3}; Assumption~\ref{as3} is used only to control \(\langle S_2,R_2\rangle_F\).
\item Theorem~\ref{thm:entrywise_threshold_support} is a standalone sparse signal-plus-noise statement.
\item Theorem~\ref{main_result_reducing_noise} uses Assumptions~\ref{as1}--\ref{as3} together with the one-sided directional stationarity hypothesis stated in the theorem; its pathwise conclusion additionally uses Definition~\ref{def:local_path_convergence}.
\item Corollary~\ref{cor:variance_scale_collapse} is a consequence of Theorem~\ref{main_result_reducing_noise}.
\item Corollary~\ref{cor:bulk_collapse_spike_stabilization} additionally uses Assumption~\ref{as3spectral}.
\item Corollary~\ref{cor:subspace_stabilization} is a further consequence of Theorem~\ref{main_result_reducing_noise} plus the asymptotic gap condition
\[
\|R_2^*(N)\|_{\op}/\delta_k(N)\to0
\qquad\text{in probability.}
\]
\item Corollary~\ref{cor:gaussian_BBP_persistent_spikes} uses Theorem~\ref{main_result_reducing_noise} together with a persistent-spike assumption \(\sigma_i^*(N)\to\bar\sigma_i>0\) and the additional exact finite-rank deformation hypothesis stated in the corollary itself.
\end{enumerate}
\end{rem}

\begin{prop}[Spectral consequence of an MP-scale admissible random component]
\label{prop:mp_bulk_interpretation}
Let \(M_N/N\to c\in(0,\infty)\). Let
\[
W_N=R_N+S_N,
\]
where \(R_N\in\R^{N\times M_N}\) has i.i.d.\ \( \mathcal N(0,g_N/N)\) entries with \(g_N\to g_\infty>0\), and \(S_N\) is independent of \(R_N\) with
\[
\rank(S_N)=o(\min\{N,M_N\}).
\]

Then:
\begin{enumerate}[leftmargin=2em]
\item the singular-value ESD \(\nu_{R_N}\), defined as in Definition~\ref{ESD_Definition} over the \(\min\{N,M_N\}\) singular values of \(R_N\), converges almost surely to the pushforward of the \emph{nonzero covariance-spectrum part} of \(\mu_{\mathrm{MP}}^{c,g_\infty}\) under \(x\mapsto \sqrt{x}\); equivalently, when \(c>1\) one first removes the MP atom at \(0\) and renormalizes. Its support is
\[
\Big[\sqrt{g_\infty}\,\big|1-\sqrt c\big|,\,\sqrt{g_\infty}(1+\sqrt c)\Big];
\]
\item the singular-value ESD of \(W_N\) has the same weak limit as \(\nu_{R_N}\), because
\[
\sup_{x\in\R}
\abs{
\nu_{W_N}((-\infty,x])-\nu_{R_N}((-\infty,x])
}
\le
\frac{\rank(S_N)}{\min\{N,M_N\}}
\to0;
\]
\item suppose, in addition, that \(S_N\) has fixed rank \(r\), its nonzero singular values \(\theta_i(N)\) converge to positive limits, and the pair \((R_N,S_N)\) satisfies the finite-rank deformation hypotheses of Theorem~\ref{thm:abstract_inverse_D} with deterministic edge \(b_+(N)\), critical threshold \(\theta_c(N)\), and inverse-\(D\) outlier map. Then every spike for which \(\PP(\theta_i(N)>\theta_c(N)+\gamma)\to1\) for some \(\gamma>0\) produces an outlier singular value \(s_i(W_N)\) satisfying
\[
s_i(W_N)-D_N^{-1}\!\left(\frac1{\theta_i(N)^2}\right)\to0
\qquad\text{in probability},
\]
and the inverse-\(D\) estimator from Theorem~\ref{thm:abstract_inverse_D} recovers \(\theta_i(N)\) in probability. In the square Gaussian case \(c=1\), this reduces to the usual BBP outlier transition.
\end{enumerate}
Thus any admissible random component that remains at MP scale leaves a persistent MP-type bulk in the singular spectrum.
\end{prop}
\proofloc{See Appendix~\ref{app:proofs_main}, Proof of Proposition~\ref{prop:mp_bulk_interpretation}.}

In this stylized additive model, Theorem~\ref{theorem_reduction_of_loss} and Proposition~\ref{prop:mp_bulk_interpretation} say that an MP-scale \(R_2\) leaves a persistent singular-value bulk, and removing it collapses that bulk while lowering the regularized objective by \(\lambda_{\mathrm{reg}}\|R_2\|_F^2\), up to the perturbation terms. Theorem~\ref{main_result_reducing_noise} further rules out admissible sub-MP Frobenius noise at locally stationary limit points, given local-path convergence.


% =================================================

\section{Proofs of the perturbation lemmas}
\label{app:proof_perturb}

\begin{lem}[Slack removal]
\label{lem:slack_removal}
Let \(X_N\) and \(Y_N\) be real random variables. Suppose that for every fixed \(\eta>0\),
\[
\big(X_N-Y_N-\eta\big)_+\to0
\qquad\text{in probability}.
\]
Then
\[
\big(X_N-Y_N\big)_+\to0
\qquad\text{in probability}.
\]
Equivalently, if \(X_N\le Y_N+\eta+o_{\PP}(1)\) for every fixed \(\eta>0\), then \(X_N\le Y_N+o_{\PP}(1)\).
\end{lem}

\begin{proof}
Fix \(\delta>0\) and set \(\eta=\delta/2\). Then
\[
\PP(X_N-Y_N>\delta)
\le
\PP(X_N-Y_N-\eta>\delta-\eta)
\le
\PP\big((X_N-Y_N-\eta)_+>\delta/2\big),
\]
and the last probability tends to \(0\) by assumption.
\end{proof}

\subsection{Proofs of the deterministic pruning certificates}

\begin{proof}[Proof of Lemma~\ref{lem:deterministic_ce_path_bound}]
For each \(s\in T\), the definition of \(L_s\) gives
\[
\|z_{A_\theta}(s)-z_{A_1}(s)\|_\infty
\le
\| \psi_2(A_\theta\psi_1(s))-\psi_2(A_1\psi_1(s))\|_\infty
\le
(1-\theta)L_s\|R\psi_1(s)\|_\infty.
\]
For the single-sample cross-entropy
\[
\ell_s(z):=-z_{C(s)}+\log\!\left(\sum_{j=1}^K e^{z_j}\right),
\]
we have \(\nabla\ell_s(z)=\rho(z)-e_{C(s)}\), hence \(\|\nabla\ell_s(z)\|_1\le2\). Therefore \(\ell_s\) is \(2\)-Lipschitz with respect to \(\|\cdot\|_\infty\). Averaging over \(s\in T\) yields
\[
\left|
\mathcal L_{\mathrm{CE}}(A_\theta)-\mathcal L_{\mathrm{CE}}(A_1)
\right|
\le
\frac2{|T|}\sum_{s\in T}(1-\theta)L_s\|R\psi_1(s)\|_\infty
=
(1-\theta)B_T(R).
\]
\end{proof}

\begin{proof}[Proof of Corollary~\ref{cor:deterministic_margin_stability}]
The logit perturbation bound in the proof of Lemma~\ref{lem:deterministic_ce_path_bound} gives
\[
\|z_{A_\theta}(s)-z_{A_1}(s)\|_\infty\le\eta_s(\theta).
\]
If the predicted label changes, then the original top logit and some competing logit must be separated by at most \(2\eta_s(\theta)\); otherwise every competing logit remains below the original top logit after perturbation. This proves the prediction-margin claim.

For accuracy, correctness is determined by the sign of the true-label margin. The true class logit can move by at most \(\eta_s(\theta)\), and the maximum competing logit can move by at most \(\eta_s(\theta)\), so the true-label margin can change by at most \(2\eta_s(\theta)\). Therefore the sign of \(\Delta_{A_1}^{\mathrm{true}}(s)\) can change only if \(|\Delta_{A_1}^{\mathrm{true}}(s)|\le2\eta_s(\theta)\). Counting such samples gives the displayed bound.
\end{proof}

\begin{proof}[Proof of Theorem~\ref{thm:deterministic_additive_frobenius_certificate}]
By Lemma~\ref{lem:deterministic_ce_path_bound},
\[
\mathcal L_{\mathrm{CE}}(A_\theta)-\mathcal L_{\mathrm{CE}}(A_1)
\le
(1-\theta)B_T(R).
\]
For the Frobenius term,
\[
\|S+\theta R\|_F^2-\|S+R\|_F^2
=
-(1-\theta^2)\|R\|_F^2
-2(1-\theta)\langle S,R\rangle_F.
\]
Multiplying this identity by \(\lambda_2\) and adding the cross-entropy bound gives the displayed path inequality. The \(\eta\)-version follows from
\[
-2\lambda_2(1-\theta)\langle S,R\rangle_F
\le
2\lambda_2(1-\theta)|\langle S,R\rangle_F|.
\]
The full-removal sufficient condition is the case \(\theta=0\).
\end{proof}

\begin{proof}[Proof of Theorem~\ref{thm:deterministic_mask_certificate}]
By Lemma~\ref{lem:deterministic_ce_path_bound},
\[
\mathcal L_{\mathrm{CE}}(A_\theta)-\mathcal L_{\mathrm{CE}}(A_1)
\le
(1-\theta)B_T(R).
\]
Since \(S\odot R=0\),
\[
\|S+\theta R\|_F^2=\|S\|_F^2+\theta^2\|R\|_F^2,
\qquad
\|S+\theta R\|_1=\|S\|_1+\theta\|R\|_1.
\]
Similarly,
\[
\|A_1\|_F^2=\|S\|_F^2+\|R\|_F^2,
\qquad
\|A_1\|_1=\|S\|_1+\|R\|_1.
\]
Substituting these identities into \(\mathcal J(A_\theta)-\mathcal J(A_1)\) gives
\[
\mathcal J(A_\theta)-\mathcal J(A_1)
\le
(1-\theta)B_T(R)
-\lambda_2(1-\theta^2)\|R\|_F^2
-\lambda_1(1-\theta)\|R\|_1.
\]
Taking \(\theta=0\) proves the full-pruning statement and the sufficient decrease condition follows immediately.
\end{proof}

\begin{proof}[Proof of Theorem~\ref{thm:deterministic_prune_restore_certificate}]
Apply Theorem~\ref{thm:deterministic_mask_certificate} to the mask decomposition
\[
A=(S+Q)+P.
\]
The pairwise disjoint-support assumption implies \((S+Q)\odot P=0\), so the theorem applies with structured part \(S+Q\) and removed part \(P\). This gives the displayed path inequality and the full prune--restore statement by taking \(\theta=0\). The sufficient decrease condition follows immediately.
\end{proof}

\begin{proof}[Proof of Corollary~\ref{cor:restoration_improves_certificate}]
The all-pruned upper bound is Theorem~\ref{thm:deterministic_mask_certificate} applied to
\[
A=S+R,
\qquad R=Q+P.
\]
The restored upper bound is Theorem~\ref{thm:deterministic_prune_restore_certificate}. Since \(Q\odot P=0\),
\[
\|R\|_F^2=\|Q\|_F^2+\|P\|_F^2,
\qquad
\|R\|_1=\|Q\|_1+\|P\|_1.
\]
Therefore the restored upper bound is smaller than the all-pruned upper bound exactly when
\[
B_T(P)-\lambda_2\|P\|_F^2-\lambda_1\|P\|_1
<
B_T(R)-\lambda_2\|R\|_F^2-\lambda_1\|R\|_1,
\]
which is equivalent to
\[
B_T(R)-B_T(P)
>
\lambda_2\|Q\|_F^2+\lambda_1\|Q\|_1.
\]
\end{proof}

\begin{proof}[Proof of Corollary~\ref{cor:weighted_mask_certificate}]
The proof is the weighted coordinatewise version of Theorem~\ref{thm:deterministic_mask_certificate}. The same cross-entropy bound applies, and \(S\odot R=0\) gives
\[
(S_{ij}+\theta R_{ij})^2=(S_{ij})^2+\theta^2(R_{ij})^2,
\qquad
|S_{ij}+\theta R_{ij}|=|S_{ij}|+\theta|R_{ij}|.
\]
Multiplying by the nonnegative weights and summing yields the displayed inequality; \(\theta=0\) gives the full-pruning condition.
\end{proof}

\begin{proof}[Proof of Corollary~\ref{cor:weighted_mask_stationarity}]
Set \(\theta=1-h\) in Corollary~\ref{cor:weighted_mask_certificate} and divide by \(h>0\). This gives
\[
\frac{\mathcal J_{\Gamma}(A_{1-h})-\mathcal J_{\Gamma}(A_1)}{h}
\le
B_T(R)
-(2-h)\sum_{i,j}\gamma_{2,ij}R_{ij}^2
-\sum_{i,j}\gamma_{1,ij}|R_{ij}|.
\]
Taking \(h\downarrow0\) and using one-sided stationarity yields
\[
-\varepsilon
\le
B_T(R)
-2\sum_{i,j}\gamma_{2,ij}R_{ij}^2
-\sum_{i,j}\gamma_{1,ij}|R_{ij}|,
\]
which is the directional claim. If the pathwise lower bound holds for every \(\theta\), combine it with Corollary~\ref{cor:weighted_mask_certificate} and divide by \(1-\theta\) to obtain the finite-\(\theta\) inequality.
\end{proof}

\begin{proof}[Proof of Theorem~\ref{thm:deterministic_mask_stationarity}]
This is Corollary~\ref{cor:weighted_mask_stationarity} with \(\gamma_{2,ij}\equiv\lambda_2\) and \(\gamma_{1,ij}\equiv\lambda_1\).
\end{proof}

\begin{proof}[Proof of Theorem~\ref{thm:deterministic_multilayer_mask_certificate}]
Apply Theorem~\ref{thm:deterministic_mask_certificate} with \(\theta=0\) at each pruning step in the network obtained after the previous \(j-1\) pruning operations. If \(\alpha^{(j)}\) denotes the parameters after step \(j\), then
\[
\mathcal J(\alpha^{(j)})-\mathcal J(\alpha^{(j-1)})
\le
-\lambda_{2,j}\|R_j\|_F^2
-\lambda_{1,j}\|R_j\|_1
+B_j(R_j).
\]
Summing over \(j=1,\dots,M\) telescopes the left-hand side and proves the displayed multi-layer bound.
\end{proof}

\begin{proof}[Proof of Corollary~\ref{cor:deterministic_multilayer_margin_stability}]
For each \(s\in T\), telescoping the logits gives
\[
\|z^{(M)}(s)-z^{(0)}(s)\|_\infty
\le
\sum_{j=1}^M\|z^{(j)}(s)-z^{(j-1)}(s)\|_\infty
\le
\eta_s^{\mathrm{tot}}.
\]
The cross-entropy bound follows from the \(2\)-Lipschitz estimate above, and the prediction and accuracy bounds follow by the margin argument of Corollary~\ref{cor:deterministic_margin_stability} with \(\eta_s(\theta)\) replaced by \(\eta_s^{\mathrm{tot}}\).
\end{proof}

\begin{proof}[Proof of Theorem~\ref{thm:deterministic_multilayer_mask_stationarity}]
Fix \(h\in(0,1]\), and let \(\alpha_h^{(j)}\) be the network after the first \(j\) fractional pruning steps, with \(\alpha_h^{(0)}=\alpha^{(0)}\) and \(\alpha_h^{(M)}=\alpha_h\). At step \(j\), apply the one-layer deterministic certificate with \(\theta=1-h\) in the intermediate network \(\alpha_h^{(j-1)}\). By the assumed validity of the fractional budget \(B_j(R_j)\),
\[
\mathcal J(\alpha_h^{(j)})-\mathcal J(\alpha_h^{(j-1)})
\le
hB_j(R_j)
-\lambda_{2,j}(2h-h^2)\|R_j\|_F^2
-\lambda_{1,j}h\|R_j\|_1.
\]
Summing over \(j=1,\dots,M\) gives
\[
\mathcal J(\alpha_h)-\mathcal J(\alpha^{(0)})
\le
h\sum_{j=1}^M B_j(R_j)
-h(2-h)\sum_{j=1}^M\lambda_{2,j}\|R_j\|_F^2
-h\sum_{j=1}^M\lambda_{1,j}\|R_j\|_1.
\]
After division by \(h\), the one-sided directional stationarity assumption and the limit \(h\downarrow0\) imply
\[
-\varepsilon
\le
\sum_{j=1}^M B_j(R_j)
-2\sum_{j=1}^M\lambda_{2,j}\|R_j\|_F^2
-\sum_{j=1}^M\lambda_{1,j}\|R_j\|_1,
\]
which is the first claim. If the stronger pathwise lower bound holds for every \(h\in[0,1]\), combine it with the displayed upper bound and divide by \(h\) to obtain the finite-\(h\) inequality.
\end{proof}

\begin{proof}[Proof of Corollary~\ref{cor:gaussian_data_path_budget}]
Condition on the vectors \(v_s=\psi_1(s)\) and on the bounds \(L_s\). For each \(s\in T\) and row \(i\),
\[
\langle r_i,v_s\rangle\sim N\!\left(0,\frac{g}{n}\|v_s\|_2^2\right).
\]
Thus
\[
\PP\!\left(
|\langle r_i,v_s\rangle|
>
\|v_s\|_2\sqrt{\frac{2g}{n}\log\frac{2n|T|}{\delta}}
\right)
\le
\frac{\delta}{n|T|}.
\]
A union bound over \(i=1,\dots,n\) and \(s\in T\) gives, with probability at least \(1-\delta\),
\[
\|Rv_s\|_\infty
\le
\|v_s\|_2\sqrt{\frac{2g}{n}\log\frac{2n|T|}{\delta}}
\qquad\forall s\in T.
\]
Substitution into \eqref{eq:data_path_budget} proves the claim.
\end{proof}

\begin{proof}[Proof of Corollary~\ref{cor:mp_edge_data_path_budget}]
For an \(n\times m\) Gaussian matrix with entries \(N(0,g/n)\), the nominal MP singular-value edge is
\[
\sigma_{+,R}=\sqrt g\left(1+\sqrt{\frac{m}{n}}\right).
\]
Equivalently,
\[
\sqrt g=\frac{\sigma_{+,R}}{1+\sqrt{m/n}}.
\]
Substituting this identity into the bound of Corollary~\ref{cor:gaussian_data_path_budget} gives the displayed inequality.
\end{proof}

\begin{proof}[Proof of Corollary~\ref{cor:multilayer_gaussian_data_path_budget}]
Apply Corollary~\ref{cor:gaussian_data_path_budget} at pruning step \(j\) with failure probability \(\delta_j\), conditional on the intermediate vectors \(v_{j,s}\) and Lipschitz bounds \(L_{j,s}\). This gives the displayed bound for \(B_j(R_j)\) with conditional probability at least \(1-\delta_j\). A union bound over \(j=1,\dots,M\) gives simultaneous validity with probability at least \(1-\sum_j\delta_j\ge1-\delta\). Substituting these simultaneous budget bounds into Theorem~\ref{thm:deterministic_multilayer_mask_certificate} proves the objective inequality.
\end{proof}

\begin{proof}[Proof of Corollary~\ref{cor:multilayer_subgaussian_data_path_budget}]
The proof is identical to Corollary~\ref{cor:multilayer_gaussian_data_path_budget}, using Corollary~\ref{cor:subgaussian_data_path_budget} in place of Corollary~\ref{cor:gaussian_data_path_budget}.
\end{proof}

\begin{proof}[Proof of Corollary~\ref{cor:subgaussian_data_path_budget}]
Condition on the variables with respect to which \(v_s\) and \(L_s\) are measurable.
Apply the assumed tail bound with
\[
t_s:=\|v_s\|_2\sqrt{\frac{g}{c n}\log\frac{2n|T|}{\delta}}.
\]
The assumed tail bound and the same row/sample union bound as above give
\[
\|Rv_s\|_\infty\le t_s
\qquad\forall s\in T
\]
with probability at least \(1-\delta\). Substitution into \eqref{eq:data_path_budget} proves the corollary.
\end{proof}

\subsection{Proof of Lemma~\ref{lem:remove_R}}

\begin{proof}
Let
\[
\nu:=\lambda(W_1s).
\]
Since \(\lambda\) is either absolute value or ReLU, it is coordinatewise \(1\)-Lipschitz and satisfies \(\|\nu\|_2\le \|W_1s\|_2\). Write
\[
X(s,\alpha_{W_2})=\lambda\big(W_3\lambda((S_2+R_2)\nu)\big),
\qquad
X(s,\alpha_{S_2})=\lambda\big(W_3\lambda(S_2\nu)\big).
\]
Using the \(1\)-Lipschitz property of \(\lambda\) and \(\|Av\|_\infty\le \|A\|_\infty\|v\|_\infty\), we obtain
\[
\big\|X(s,\alpha_{W_2})-X(s,\alpha_{S_2})\big\|_\infty
\le
\|W_3\|_\infty\,\|R_2\nu\|_\infty.
\]
So it suffices to bound \(\|R_2\nu\|_\infty\).

By Assumption~\ref{as2}, \(R_2\) is independent of \(W_1\), hence independent of \(\nu\). Conditioned on \(\nu\), each coordinate \((R_2\nu)_j\) is centered Gaussian with variance
\[
\mathrm{Var}\big((R_2\nu)_j\mid \nu\big)
=
\frac{g(t)}{N}\|\nu\|_2^2
\le
\frac{\|W_1s\|_2^2}{N}.
\]
Set
\[
u_N:=2\sqrt{2}\,\|W_1s\|_2\sqrt{\frac{\log(2N)}{N}}.
\]
If \(\|W_1s\|_2=0\), then \(\nu=0\) and the claim is immediate. Otherwise, the elementary Gaussian tail bound gives, conditionally on \(\nu\),
\[
\PP\!\left(|(R_2\nu)_j|>u_N\,\middle|\,\nu\right)
\le
2\exp\!\left(
-\frac{u_N^2}{2\|W_1s\|_2^2/N}
\right)
=2(2N)^{-4}.
\]
A union bound over \(j=1,\dots,N\) gives
\[
\PP\!\left(
\|R_2\nu\|_\infty >
\;2\sqrt{2}\,\|W_1s\|_2\sqrt{\frac{\log(2N)}{N}}
\middle|\,\nu
\right)
\le
\frac{1}{8N^3}
\le
\frac{1}{N}.
\]
The same bound holds unconditionally after averaging over \(\nu\).
Multiplying by \(\|W_3\|_\infty\) yields the desired estimate.
\end{proof}

\subsection{Proof of Lemma~\ref{lem:remove_R_general}}

\begin{proof}
Let \(u:=\psi_1(s)\). Since
\[
X(s,\alpha_W)=\psi_2((R+S)u),
\qquad
X(s,\alpha_S)=\psi_2(Su),
\]
Assumption~\ref{as4} gives
\[
\|X(s,\alpha_W)-X(s,\alpha_S)\|_\infty
\le
L_{\psi_2}(s)\,\|Ru\|_\infty.
\]
Applying Assumption~\ref{as5} with \(v=u\) and then averaging over the conditioning variables gives
\[
\PP\big(\|Ru\|_\infty>d_1(N)J(u)\big)\le d_2(N).
\]
Therefore
\[
\PP\!\left(
\|X(s,\alpha_W)-X(s,\alpha_S)\|_\infty
\le
d_1(N)J(\psi_1(s))\,L_{\psi_2}(s)
\right)\ge 1-d_2(N),
\]
which is the claim.
\end{proof}

\subsection{Justification of Example~\ref{ex:subgaussian_rows}}

The standard sub-Gaussian tail bound gives absolute constants \(c,C>0\) such that, conditionally on \(S,\psi_1,\psi_2\),
\[
\PP\!\left(
|\langle r_j,v\rangle|>u
\,\middle|\,S,\psi_1,\psi_2
\right)
\le
2\exp\!\left(-c\frac{Nu^2}{\kappa^2\|v\|_2^2}\right).
\]
Taking
\[
u=C\kappa\|v\|_2\sqrt{\frac{\log(2N/\delta)}{N}}
\]
with \(C\) large enough and applying a union bound over \(j=1,\dots,N\) yields the displayed \(\ell_\infty\) concentration bound.

\subsection{Proof of Lemma~\ref{lem:remove_R_loosandacc}}

\begin{proof}
For \(s\in T\), let \(X(\alpha,s)\in\R^K\) denote the vector of pre-softmax logits. By Lemma~\ref{lem:remove_R}, for each fixed \(s\in T\),
\[
\PP\Big(
\|X(\alpha_{S_2},s)-X(\alpha_{W_2},s)\|_\infty\le \tau_N(s)
\Big)
\ge
1-\frac{1}{N},
\]
where
\[
\tau_N(s)
=
\;2\sqrt{2}\,\|W_3\|_\infty\|W_1s\|_2\sqrt{\frac{\log(2N)}{N}}
=2\sqrt{2}\,a(N,s).
\]
Let \(\ell_s(z)=-z_{C(s)}+\log\sum_j e^{z_j}\). As in the proof of Lemma~\ref{lem:deterministic_ce_path_bound}, \(\ell_s\) is \(2\)-Lipschitz with respect to \(\|\cdot\|_\infty\).
Applying this with \(z=X(\alpha_{S_2},s)\) and \(z'=X(\alpha_{W_2},s)\), we get
\[
|\ell_s(X(\alpha_{S_2},s))-\ell_s(X(\alpha_{W_2},s))|
\le
4\sqrt{2}\,a(N,s)
\]
on the high-probability event from Lemma~\ref{lem:remove_R}. Since \(T\) is finite, a union bound over \(s\in T\) yields
\[
\PP\!\left(
|L_{\cancel{\mathrm{reg}}}(\alpha_{S_2})-L_{\cancel{\mathrm{reg}}}(\alpha_{W_2})|
\le
\frac{4\sqrt{2}}{|T|}\sum_{s\in T}a(N,s)
\right)
\ge
1-\frac{|T|}{N},
\]
which proves the lemma.
\end{proof}

\subsection{Proof of Theorem~\ref{thm:finite_confidence_gaussian}}

\begin{proof}
For each \(s\in T_0\), set \(\nu_s=\lambda(W_1s)\). Applying the same conditional Gaussian row bound and union bound as in Lemma~\ref{lem:remove_R}, now over \(s\in T_0\) and \(j=1,\dots,N\) with failure probability \(\delta\), gives the displayed simultaneous bound on \(\|R_2\nu_s\|_\infty\). Multiplying by \(\|W_3\|_\infty\) yields the logit bound, and the cross-entropy bound follows from the \(2\)-Lipschitz estimate for \(\ell_s\).

For the regularized statement, use the exact identity
\[
L(\alpha_{S_2})-L(\alpha_{W_2})
=
L_{\cancel{\mathrm{reg}}}(\alpha_{S_2})-L_{\cancel{\mathrm{reg}}}(\alpha_{W_2})
-\lambda_{\mathrm{reg}}\|R_2\|_F^2
-2\lambda_{\mathrm{reg}}\langle S_2,R_2\rangle_F.
\]
Conditional on \(S_2\),
\[
\langle S_2,R_2\rangle_F
\sim
\mathcal N\!\left(0,\frac{g(t)}{N}\|S_2\|_F^2\right),
\]
and Assumption~\ref{as3} gives
\[
\PP\big(|\langle S_2,R_2\rangle_F|>\eta\mid S_2\big)
\le
2e^{-N\eta^2/(2B_N^2)}.
\]
Intersecting the logit event with \(\{|\langle S_2,R_2\rangle_F|\le\eta\}\) gives the claimed probability and inequality.
\end{proof}



\subsection{Proof of Theorem~\ref{thm:entrywise_threshold_support}}

\begin{proof}
Condition on \(S_N\) when it is random; the remaining probability is over \(R_N\).
If \(g_N=0\), then \(R_N=0\) almost surely and the claim is immediate. Assume \(g_N>0\). 
For every entry,
\[
\PP(|R_{ij}|>t_{N,\delta})
\le
2\exp\!\left(-\frac{Nt_{N,\delta}^2}{2g_N}\right)
=
\frac{\delta}{NM_N},
\]
and a union bound gives
\[
\PP\!\left(\max_{i,j}|R_{ij}|\le t_{N,\delta}\right)\ge1-\delta.
\]
On this event, if \((i,j)\notin\operatorname{supp}(S_N)\), then \(|W_{ij}|=|R_{ij}|\le t_{N,\delta}\), so \(\mathcal H_{t_{N,\delta}}\) removes that entry. Thus thresholding creates no false positives. If \(|(S_N)_{ij}|>2t_{N,\delta}\), then
\[
|W_{ij}|\ge |(S_N)_{ij}|-|R_{ij}|>2t_{N,\delta}-t_{N,\delta}=t_{N,\delta},
\]
so thresholding keeps that entry. This proves the support sandwich and the missed-support inclusion. On the beta-min event \(\mathcal B_{N,\delta}\), every element of \(\operatorname{supp}(S_N)\) has magnitude larger than \(2t_{N,\delta}\), so the support sandwich collapses to exact support recovery. Unconditionally, the exact-recovery failure probability is at most \(\delta+\PP(\mathcal B_{N,\delta}^c)\le\delta+\gamma\).
\end{proof}

\section{Low-rank deformations and proofs of the main theorems}
\label{app:proofs_main}

\subsection{External random-matrix inputs quoted in the manuscript}
\label{app:external_rmt_inputs}

The Marchenko--Pastur theorem in Theorem~\ref{RMT_MP_theorem} and the square Gaussian finite-rank deformation theorem in Theorem~\ref{thm:BGN_square} are used as standard random-matrix inputs. The statements are classical; see, for example, \cite{marchenko1967distribution} for Marchenko--Pastur theory and \cite{benaych2012singular} for the Benaych-Georges--Nadakuditi outlier theorem. All paper-specific consequences stated after these inputs are proved below.

\subsection{Finite-rank deformations of random matrices}

\begin{thm}[Benaych-Georges--Nadakuditi, square Gaussian case \cite{benaych2012singular}]
\label{thm:BGN_square}
Let \(W_N=R_N+S_N\), where \(R_N\) is an \(N\times N\) random matrix with i.i.d.\ \(N(0,1/N)\) entries and \(S_N\) is deterministic, or independent of \(R_N\), with fixed rank \(r\) and nonzero singular values \(\sigma_1>\cdots>\sigma_r>0\) independent of \(N\). Then the \(i\)-th largest singular value \(s_i(W_N)\) converges almost surely to
\[
s_i(W_N)\to
\begin{cases}
\sigma_i+\sigma_i^{-1}, & \sigma_i>1,\\[0.2em]
2, & \sigma_i\le 1.
\end{cases}
\]
In particular, subcritical spikes merge with the bulk edge \(2\).
\end{thm}
\proofloc{This is a classical cited input; see the source note in Appendix~\ref{app:external_rmt_inputs}. The variance-\(g\) and varying-scale consequences are proved immediately below.}

\begin{cor}[Square Gaussian case with variance \(g/N\)]
\label{cor:BGN_square_g}
Let \(W_N=R_N+S_N\), where \(R_N\) is an \(N\times N\) random matrix with i.i.d.\ \(N(0,g/N)\) entries for some fixed \(g>0\), and \(S_N\) is deterministic, or independent of \(R_N\), with fixed rank \(r\) and distinct nonzero singular values \(\sigma_1>\cdots>\sigma_r>0\). Then
\[
s_i(W_N)\to
\begin{cases}
\sigma_i+g/\sigma_i, & \sigma_i>\sqrt g,\\[0.2em]
2\sqrt g, & \sigma_i\le \sqrt g.
\end{cases}
\]
\end{cor}

\begin{proof}
Write
\[
W_N=\sqrt g\,(\widetilde R_N+\widetilde S_N),
\qquad
\widetilde R_N:=R_N/\sqrt g,
\qquad
\widetilde S_N:=S_N/\sqrt g.
\]
Then \(\widetilde R_N\) has i.i.d.\ \(N(0,1/N)\) entries, and the nonzero singular values of \(\widetilde S_N\) are \(\sigma_i/\sqrt g\). Apply Theorem~\ref{thm:BGN_square} and multiply by \(\sqrt g\).
\end{proof}

\begin{thm}[Square Gaussian outlier map with varying spike and noise scales]
\label{thm:BGN_square_variable}
Let \(W_N=R_N+S_N\), where \(R_N\) is an \(N\times N\) random matrix with i.i.d.\ \(N(0,g_N/N)\) entries and \(g_N\to g_\infty\in[0,\infty)\). Assume \(S_N\) has fixed rank \(r\). Write its nonzero singular values in decreasing order,
\[
\sigma_1(S_N)\ge\cdots\ge\sigma_r(S_N)>0,
\]
and assume that for each fixed \(1\le i\le r\),
\[
\sigma_i(S_N)\to \bar\sigma_i\in(0,\infty)
\qquad\text{in probability}
\]
where \(\bar\sigma_1\ge\cdots\ge\bar\sigma_r>0\), and that \(S_N\) is independent of \(R_N\).
Then for each fixed \(1\le i\le r\),
\[
s_i(W_N)\to
\begin{cases}
\bar\sigma_i+g_\infty/\bar\sigma_i, & \bar\sigma_i>\sqrt{g_\infty},\\[0.2em]
2\sqrt{g_\infty}, & \bar\sigma_i\le \sqrt{g_\infty}.
\end{cases}
\qquad\text{in probability.}
\]
Moreover, when \(g_\infty>0\), the usual BGN alignment conclusions hold for simple separated spikes: supercritical spikes have asymptotically nonzero projection onto the corresponding signal directions, while subcritical spikes have vanishing projection. For repeated or clustered spikes, these statements are understood as subspace-alignment statements. When \(g_\infty=0\), any separated spike cluster has convergent spectral projectors under the Wedin gap condition \(\|R_N\|_{\op}/\mathrm{gap}_N\to0\).
\end{thm}

\begin{proof}
Condition on \(S_N\). First suppose \(g_\infty>0\). Then \(g_N>0\) for all large \(N\), and after dividing by \(\sqrt{g_N}\) one reduces to the unit-variance square Gaussian finite-rank deformation theorem of Benaych-Georges and Nadakuditi \cite{benaych2012singular}. On events whose probability tends to one, the conditioned spike values are arbitrarily close to their deterministic limits; the BGN outlier map is continuous away from the threshold, and at the threshold both branches give the edge. Thus the conditional deterministic-spike result and Slutsky's theorem give the displayed convergence in probability. The limiting outlier map in the unit-variance scaling is
\[
\theta\mapsto \theta+\theta^{-1},
\]
so in the original scaling it becomes
\[
\sigma\mapsto \sigma+\frac{g_\infty}{\sigma},
\]
with critical threshold \(\sqrt{g_\infty}\) and edge \(2\sqrt{g_\infty}\). The singular-vector and spike-subspace alignment statements in the positive-noise case are exactly the corresponding supercritical/subcritical BGN conclusions.

Now suppose \(g_\infty=0\). Write \(R_N=\sqrt{g_N}\,Z_N\), where \(Z_N\) has i.i.d.\ \(N(0,1/N)\) entries. By the standard Gaussian operator-norm bound, \(\|Z_N\|_{\op}=O_{\PP}(1)\), hence
\[
\|R_N\|_{\op}\to0
\qquad\text{in probability}.
\]
Weyl's singular-value inequality gives, for each fixed \(i\),
\[
|s_i(W_N)-s_i(S_N)|\le \|R_N\|_{\op}\to0
\qquad\text{in probability}.
\]
Since \(s_i(S_N)\to\bar\sigma_i\), we obtain \(s_i(W_N)\to\bar\sigma_i\), which is the asserted formula with \(g_\infty=0\). The stated projector convergence in the separated-cluster case follows from Wedin's theorem.
\end{proof}

\begin{thm}[Abstract inverse-\(D\) spike recovery]
\label{thm:abstract_inverse_D}
Let \(W_N=R_N+S_N\). Assume that for each fixed \(1\le i\le r\), the \(i\)-th nonzero singular value of \(S_N\) is \(\theta_i(N)>0\). Assume further that there exist deterministic quantities \(b_+(N)\), \(\theta_c(N)\), and a strictly decreasing map
\[
D_N:(b_+(N),\infty)\to (0,D_N(b_+(N)^+))
\]
with the following properties:
\begin{enumerate}[leftmargin=2em]
\item on supercritical spike sequences, meaning \(\PP(\theta_i(N)>\theta_c(N))\to1\), one has
\[
s_i(W_N)-\rho_N(\theta_i(N))\to0
\qquad\text{in probability},
\]
where
\[
\rho_N(\theta):=D_N^{-1}\!\left(\frac1{\theta^2}\right);
\]
\item on subcritical spike sequences, meaning \(\PP(\theta_i(N)\le \theta_c(N))\to1\), one has
\[
s_i(W_N)-b_+(N)\to0
\qquad\text{in probability};
\]
\item in the supercritical case, the singular vectors associated with \(s_i(W_N)\) have asymptotically nonvanishing projection onto the corresponding singular directions of \(S_N\);
\item for each fixed supercritical spike under consideration, \(\PP(\theta_i(N)>\theta_c(N))\to1\), \(\theta_i(N)\to\bar\theta_i\in(0,\infty)\), and there exist deterministic intervals
\[
I_{i,N}\subset (b_+(N),\infty)
\]
and deterministic constants \(L_{i,N}<\infty\) such that \(D_N\) is \(L_{i,N}\)-Lipschitz on \(I_{i,N}\) and
\[
\begin{aligned}
&\PP\big(\rho_N(\theta_i(N))\in I_{i,N}\big)\to1,
\qquad
\PP\big(s_i(W_N)\in I_{i,N}\big)\to1,\\
&L_{i,N}\big|s_i(W_N)-\rho_N(\theta_i(N))\big|\to0
\quad\text{in probability.}
\end{aligned}
\]
\end{enumerate}
Define the inverse-\(D\) estimator on the event \(s_i(W_N)>b_+(N)\), whose probability tends to one for every fixed supercritical spike under consideration, by
\[
\widehat\theta_i^{(D)}(N):=D_N(s_i(W_N))^{-1/2},
\]
and define it arbitrarily off this event. Then for every fixed supercritical spike,
\[
\widehat\theta_i^{(D)}(N)-\theta_i(N)\to0
\qquad\text{in probability}.
\]
\end{thm}

\begin{proof}
On the event \(\{s_i(W_N)\in I_{i,N},\ \rho_N(\theta_i(N))\in I_{i,N}\}\), whose probability tends to one, the Lipschitz assumption gives
\[
\big|D_N(s_i(W_N))-D_N(\rho_N(\theta_i(N)))\big|
\le
L_{i,N}\big|s_i(W_N)-\rho_N(\theta_i(N))\big|.
\]
Therefore
\[
D_N(s_i(W_N))
=
D_N(\rho_N(\theta_i(N)))+o_{\PP}(1)
=
\frac1{\theta_i(N)^2}+o_{\PP}(1).
\]
Since \(\theta_i(N)\to\bar\theta_i\in(0,\infty)\), the quantity \(\theta_i(N)^{-2}\) stays in a compact subinterval of \((0,\infty)\) for all large \(N\). Applying the continuous map \(x\mapsto x^{-1/2}\) on that compact interval yields
\[
\widehat\theta_i^{(D)}(N)-\theta_i(N)\to0
\qquad\text{in probability.}
\]
This is the inverse-outlier-map viewpoint underlying \(D\)-transform shrinkage methods such as OptShrink \cite{nadakuditi2014optshrink}.
\end{proof}

\subsection{Proof of Proposition~\ref{prop:mp_bulk_interpretation}}

\begin{proof}
Part (i) is the MP law on the covariance scale together with the pushforward \(x\mapsto \sqrt x\) to the singular-value scale. Because \(\nu_{R_N}\) averages over only \(\min\{N,M_N\}\) singular values, when \(c>1\) one removes the covariance-side zero atom of the full MP law before pushing forward and renormalizes the remaining mass. Part (ii) follows from the rank inequality for empirical spectral distributions: if \(A,B\) are matrices of the same dimensions, then
\[
\sup_x \abs{\nu_A((-\infty,x])-\nu_B((-\infty,x])}
\le
\frac{\rank(A-B)}{\min\{N,M_N\}}.
\]
Applying this with \(A=W_N\), \(B=R_N\), and \(A-B=S_N\) gives the claimed bound. Part (iii) is the rectangular outlier principle encoded by the inverse-\(D\) transform in Theorem~\ref{thm:abstract_inverse_D}; the square Gaussian BBP statement follows as the special case \(c=1\).
\end{proof}

\subsection{Proofs of Theorem~\ref{theorem_reduction_of_loss}, Corollary~\ref{cor:scaled_reduction_of_loss}, Corollary~\ref{cor:margin_accuracy}, and Theorem~\ref{main_result_reducing_noise} (Gaussian instances)}

\begin{proof}[Proof of Theorem~\ref{theorem_reduction_of_loss}]
Conditional on \(S_2(t)\), Assumption~\ref{as2} gives
\[
\langle S_2(t),R_2(t)\rangle_F
=
\tr(S_2(t)^\top R_2(t))
\sim
\mathcal N\!\left(0,\frac{g(t)}{N}\|S_2(t)\|_F^2\right).
\]
Since \(\|S_2(t)\|_F\le B_N\) and \(g(t)\le1\), for every \(\eta>0\),
\begin{equation}
\label{eq:cross-term-concentration-app}
\PP\!\left(
|\langle S_2(t),R_2(t)\rangle_F|>\eta
\,\middle|\,S_2(t)
\right)
\le
2e^{-N\eta^2/(2B_N^2)}.
\end{equation}
Hence
\begin{equation}
\label{eq:cross-term-oP-app}
\langle S_2(t),R_2(t)\rangle_F\to 0
\qquad\text{in probability as }N\to\infty.
\end{equation}

Write
\[
L(\alpha)=L_{\cancel{\mathrm{reg}}}(\alpha)+\lambda_{\mathrm{reg}}\sum_{i=1}^L \|W_i\|_F^2.
\]
If \(W_2=R_2+S_2\) and we replace \(W_2\) by \(S_2\), then
\begin{align}
L(\alpha_{S_2})-L(\alpha_{W_2})
&=
\big(L_{\cancel{\mathrm{reg}}}(\alpha_{S_2})-L_{\cancel{\mathrm{reg}}}(\alpha_{W_2})\big)
-
\lambda_{\mathrm{reg}}\big(\|W_2\|_F^2-\|S_2\|_F^2\big)\notag\\
&=
\big(L_{\cancel{\mathrm{reg}}}(\alpha_{S_2})-L_{\cancel{\mathrm{reg}}}(\alpha_{W_2})\big)
-
\lambda_{\mathrm{reg}}\|R_2\|_F^2
-
2\lambda_{\mathrm{reg}}\langle S_2,R_2\rangle_F.
\label{eq:loss-diff-exact-app}
\end{align}
By Lemma~\ref{lem:remove_R_loosandacc},
\[
\PP\!\left(
|L_{\cancel{\mathrm{reg}}}(\alpha_{S_2})-L_{\cancel{\mathrm{reg}}}(\alpha_{W_2})|
\le
\frac{4\sqrt{2}}{|T|}\sum_{s\in T}a(N,s)
\right)
\ge
1-\frac{|T|}{N}.
\]
Intersecting this event with \(|\langle S_2,R_2\rangle_F|\le \eta\), whose probability is bounded below by \(1-2e^{-N\eta^2/(2B_N^2)}\), yields the finite-\(N\) inequality in the theorem. Letting \(N\to\infty\) and using \eqref{eq:cross-term-oP-app} gives the asymptotic expansion.
\end{proof}

\begin{proof}[Proof of Corollary~\ref{cor:scaled_reduction_of_loss}]
Let \(W_2(\vartheta)=S_2+\vartheta R_2\). Then
\begin{align}
L(\alpha_{W_2(\vartheta)})-L(\alpha_{W_2})
&=
\big(L_{\cancel{\mathrm{reg}}}(\alpha_{W_2(\vartheta)})-L_{\cancel{\mathrm{reg}}}(\alpha_{W_2})\big)\notag\\
&\quad-
\lambda_{\mathrm{reg}}
\Big(
(1-\vartheta^2)\|R_2\|_F^2
+
2(1-\vartheta)\langle S_2,R_2\rangle_F
\Big).
\label{eq:scaled-loss-diff-app}
\end{align}
For each \(s\in T\), write \(\nu_s=\lambda(W_1s)\). Let \(E_N^{\mathrm{logit}}\) be the event on which
\[
\|W_3\|_\infty\|R_2\nu_s\|_\infty\le \tau_N(s)
\qquad\forall s\in T.
\]
By the same samplewise Gaussian estimate as in Lemma~\ref{lem:remove_R} and a union bound over \(T\),
\[
\PP(E_N^{\mathrm{logit}})\ge 1-\frac{|T|}{N}.
\]
On \(E_N^{\mathrm{logit}}\), for every \(\vartheta\in[0,1]\) and every \(s\in T\),
\[
\big\|X(s,\alpha_{W_2(\vartheta)})-X(s,\alpha_{W_2})\big\|_\infty
\le
(1-\vartheta)\|W_3\|_\infty \|R_2\nu_s\|_\infty
\le
(1-\vartheta)\tau_N(s).
\]
The \(2\)-Lipschitz property of the cross-entropy therefore gives, on the same event,
\[
\big|L_{\cancel{\mathrm{reg}}}(\alpha_{W_2(\vartheta)})-L_{\cancel{\mathrm{reg}}}(\alpha_{W_2})\big|
\le
\frac{4\sqrt{2}(1-\vartheta)}{|T|}\sum_{s\in T}a(N,s)
\]
for every \(\vartheta\in[0,1]\). Combining \(E_N^{\mathrm{logit}}\) with \(|\langle S_2,R_2\rangle_F|\le \eta\) proves the finite-\(N\) bound. The asymptotic expansion follows from \eqref{eq:scaled-loss-diff-app} and \eqref{eq:cross-term-oP-app}. Since the finite-\(N\) event is uniform in \(\vartheta\), the same expansion also holds for deterministic sequences \(\vartheta_N\in[0,1]\).
\end{proof}

\begin{proof}[Proof of Corollary~\ref{cor:margin_accuracy}]
By Lemma~\ref{lem:remove_R}, for each fixed \(s\in T'\),
\[
\PP\!\left(
\|X(s,\alpha_{W_2})-X(s,\alpha_{S_2})\|_\infty\le \tau_N(s)
\right)\ge 1-\frac{1}{N}.
\]
A union bound over \(s\in T'\) gives simultaneous vector control with probability at least \(1-\frac{|T'|}{N}\). On that event,
\[
|\delta X(s,\alpha_{W_2})-\delta X(s,\alpha_{S_2})|\le 2\tau_N(s)
\qquad\forall s\in T'.
\]
Therefore labels can only change on samples with \(|\delta X(s,\alpha_{W_2})|\le 2\tau_N(s)\), which yields the claim.
\end{proof}

\begin{proof}[Proof of Theorem~\ref{main_result_reducing_noise}]
Set
\[
b_N:=\frac{4\sqrt{2}}{|T|}\sum_{s\in T}a_*(N,s)\to 0.
\]
Choose a deterministic sequence \(\eta_N\downarrow0\) such that
\[
\frac{N\eta_N^2}{B_N^2}\to\infty;
\]
such a sequence exists because \(B_N=o(\sqrt N)\). By Corollary~\ref{cor:scaled_reduction_of_loss} applied at the limit point, there exists an event \(E_N\) with \(\PP(E_N)\to1\) such that on \(E_N\), for every \(\vartheta\in[0,1]\),
\[
L(\alpha_\vartheta^*(N))-L(\alpha^*(N))
\le
-\lambda_{\mathrm{reg}}(1-\vartheta^2)\|R_2^*(N)\|_F^2
+(1-\vartheta)b_N
+2\lambda_{\mathrm{reg}}(1-\vartheta)\eta_N
\]
Write \(\vartheta=1-h\). Dividing by \(h>0\) gives
\[
\frac{L(\alpha_{1-h}^*(N))-L(\alpha^*(N))}{h}
\le
-\lambda_{\mathrm{reg}}(2-h)\|R_2^*(N)\|_F^2
+b_N
+2\lambda_{\mathrm{reg}}\eta_N
\]
for every \(h\in(0,1]\). On the high-probability event where the one-sided directional stationarity hypothesis also holds, taking \(\liminf_{h\downarrow0}\) yields
\[
2\lambda_{\mathrm{reg}}\|R_2^*(N)\|_F^2
\le
b_N+\epsilon_N+2\lambda_{\mathrm{reg}}\eta_N.
\]
Therefore
\[
\|R_2^*(N)\|_F^2
\le
\frac{b_N+\epsilon_N}{2\lambda_{\mathrm{reg}}}+\eta_N
+o_{\PP}(1).
\]
Since \(\eta_N\to0\), this proves the quantitative estimate in the theorem. The vanishing-in-probability statement follows immediately when \(b_N+\epsilon_N\to0\) in probability.

Now assume the local-path convergence hypothesis from Definition~\ref{def:local_path_convergence}. Fix \(\epsilon>0\). For each fixed \(N\),
\[
\|R_2(t,N)\|_F\to \|R_2^*(N)\|_F
\qquad\text{in probability as }t\to\infty
\]
by the continuous mapping theorem. Since \(R_2^*(N)\) is Gaussian (possibly degenerate if \(g_*(N)=0\)), the random variable \(\|R_2^*(N)\|_F\) has no atom at any positive \(\epsilon\). Hence
\[
\PP\big(\|R_2(t,N)\|_F\le \epsilon\big)\to \PP\big(\|R_2^*(N)\|_F\le \epsilon\big)
\qquad\text{as }t\to\infty.
\]
Letting \(N\to\infty\) and using the limit-point statement proves the theorem.
\end{proof}

\begin{proof}[Proof of Corollary~\ref{cor:variance_scale_collapse}]
Suppose, toward a contradiction, that there exist \(\eta_0>0\) and a subsequence such that \(g_*(N)N\ge \eta_0\). Since
\[
\|R_2^*(N)\|_F^2
=
\frac{g_*(N)}{N}\sum_{i,j=1}^{N} Z_{ij}^2,
\qquad
Z_{ij}\overset{\mathrm{i.i.d.}}{\sim}N(0,1),
\]
we have
\[
\frac{\|R_2^*(N)\|_F^2}{g_*(N)N}
=
\frac{1}{N^2}\sum_{i,j=1}^{N}Z_{ij}^2
\to 1
\qquad\text{in probability}.
\]
Therefore
\[
\PP\!\left(\|R_2^*(N)\|_F^2\ge \eta_0/2\right)\to 1,
\]
contradicting Theorem~\ref{main_result_reducing_noise}. Hence \(g_*(N)N\to0\). The edge statement follows from the square-case MP singular-value edge \(2\sqrt{g_*(N)}\).
In the rectangular \(N\times M_N\) case,
\[
\|R_2^*(N)\|_F^2
=
\frac{g_*(N)}{N}\sum_{i=1}^{N}\sum_{j=1}^{M_N}Z_{ij}^2,
\qquad
\frac{1}{NM_N}\sum_{i,j}Z_{ij}^2\to1
\]
in probability. Thus \(\|R_2^*(N)\|_F\to0\) forces \(g_*(N)M_N\to0\), and if \(M_N/N\to c\in(0,\infty)\) this is equivalent to \(g_*(N)N\to0\).
\end{proof}

\begin{proof}[Proof of Corollary~\ref{cor:bulk_collapse_spike_stabilization}]
Since \(\|R_2^*(N)\|_{\op}\le \|R_2^*(N)\|_F\), Theorem~\ref{main_result_reducing_noise} immediately gives
\[
\|R_2^*(N)\|_{\op}\to 0
\qquad\text{in probability}.
\]
This proves (i).

For (ii), singular-value Weyl perturbation gives
\[
\big|s_i(W_2^*(N))-s_i(S_2^*(N))\big|
\le
\|R_2^*(N)\|_{\op}
\qquad\forall i.
\]
Hence for each fixed \(i\ge1\),
\[
\big|s_i(W_2^*(N))-\sigma_i^*(N)\big|
\le
\|R_2^*(N)\|_{\op}\to0
\]
in probability.

Let \(S_{2,\le k_N}^*(N)\) denote the best rank-\(k_N\) approximation of \(S_2^*(N)\). By the Eckart--Young theorem,
\[
\sum_{i>k_N} s_i(W_2^*(N))^2
=
\inf_{\rank(A)\le k_N}\|W_2^*(N)-A\|_F^2
\le
\|W_2^*(N)-S_{2,\le k_N}^*(N)\|_F^2
\]
\[
=
\|R_2^*(N)+(S_2^*(N)-S_{2,\le k_N}^*(N))\|_F^2
\le
2\|R_2^*(N)\|_F^2+2\sum_{i>k_N}\sigma_i^*(N)^2,
\]
which proves (iii) using Theorem~\ref{main_result_reducing_noise} and Assumption~\ref{as3spectral}.

For (iv), fix \(\epsilon>0\). Since \(\|W_2^*(N)\|_F\le \|S_2^*(N)\|_F+\|R_2^*(N)\|_F\), we have
\[
\nu_{W_2^*(N)}((\epsilon,\infty))
\le
\frac{\|W_2^*(N)\|_F^2}{N\epsilon^2}
\le
\frac{2B_N^2+2\|R_2^*(N)\|_F^2}{N\epsilon^2},
\]
which converges to \(0\) in probability because \(B_N=o(\sqrt N)\) and \(\|R_2^*(N)\|_F\to0\) in probability. The final claim follows from (ii).
\end{proof}

\begin{proof}[Proof of Corollary~\ref{cor:subspace_stabilization}]
By Corollary~\ref{cor:bulk_collapse_spike_stabilization}(i),
\[
\|R_2^*(N)\|_{\op}\to0
\qquad\text{in probability.}
\]
Hence the event
\[
F_N:=\left\{\|R_2^*(N)\|_{\op}\le \frac{\delta_k(N)}{4}\right\}
\]
satisfies \(\PP(F_N)\to1\).

On \(F_N\), Weyl's singular-value inequality gives
\[
s_k(W_2^*(N))
\ge
\sigma_k^*(N)-\|R_2^*(N)\|_{\op},
\]
while
\[
s_{k+1}(W_2^*(N))
\le
\sigma_{k+1}^*(N)+\|R_2^*(N)\|_{\op},
\]
with the convention \(\sigma_j^*(N)=0\) beyond the rank of \(S_2^*(N)\). Therefore,
\[
s_k(W_2^*(N))-s_{k+1}(W_2^*(N))
\ge
\delta_k(N)-2\|R_2^*(N)\|_{\op}
\ge
\frac{\delta_k(N)}{2}
\]
on \(F_N\). Thus on \(F_N\), the top-\(k\) singular subspace of \(W_2^*(N)\) is separated from the rest by a positive gap.

Applying Wedin's \(\sin\Theta\) theorem for singular subspaces \cite{wedin1972perturbation,stewart1990matrix} to
\[
W_2^*(N)=S_2^*(N)+R_2^*(N),
\]
with target rank \(k\), yields a universal constant \(C>0\) such that on \(F_N\),
\[
\|\widehat P_{U,k}(N)-P_{U_*,k}(N)\|_{\op}
\le
C\frac{\|R_2^*(N)\|_{\op}}{\delta_k(N)},
\qquad
\|\widehat P_{V,k}(N)-P_{V_*,k}(N)\|_{\op}
\le
C\frac{\|R_2^*(N)\|_{\op}}{\delta_k(N)}.
\]
Since \(\|R_2^*(N)\|_{\op}/\delta_k(N)\to0\) in probability and \(\PP(F_N)\to1\), both projector differences converge to \(0\) in probability.
\end{proof}

\begin{proof}[Proof of Corollary~\ref{cor:gaussian_BBP_persistent_spikes}]
By Corollary~\ref{cor:variance_scale_collapse},
\[
g_*(N)N\to0,
\]
hence \(g_*(N)\to0\). Since \(\sigma_i^*(N)\to\bar\sigma_i>0\) in probability for each \(1\le i\le k\),
\[
\PP\!\big(\sigma_i^*(N)>\sqrt{g_*(N)}\big)\to1,
\]
which proves (i).

By Corollary~\ref{cor:bulk_collapse_spike_stabilization}(i),
\[
\|R_2^*(N)\|_{\op}\to0
\qquad\text{in probability}.
\]
Hence Weyl's inequality gives
\[
s_i(W_2^*(N))-\sigma_i^*(N)\to0
\qquad\text{in probability}
\]
for each fixed \(1\le i\le k\). Since \(\sigma_i^*(N)\to\bar\sigma_i>0\) and \(2\sqrt{g_*(N)}\to0\), these singular values lie above the Gaussian bulk edge \(2\sqrt{g_*(N)}\) with probability tending to one. Equivalently, in the exact finite-rank Gaussian deformation language of Theorem~\ref{thm:BGN_square_variable}, the persistent spikes are eventually supercritical. If the additional gap condition in (ii) holds for a spike, Wedin's theorem applied to \(W_2^*(N)=S_2^*(N)+R_2^*(N)\) gives, for the associated left and right rank-one projectors,
\[
\|\widehat P_i(N)-P_i^*(N)\|_{\op}
\le
C\frac{\|R_2^*(N)\|_{\op}}{\Delta_i(N)}
\to0
\qquad\text{in probability},
\]
and similarly for the right projector. This proves the strengthened projector statement in (ii).

Set \(y_N=s_i(W_2^*(N))\). The preceding Weyl estimate gives \(y_N-\sigma_i^*(N)\to0\) in probability, hence \(y_N\to\bar\sigma_i>0\). Also \(g_*(N)\to0\), so the event \(y_N^2>4g_*(N)\) has probability tending to one. On this event,
\[
0\le
y_N-\widehat\sigma_i^{\mathrm{BBP}}(N)
=
\frac12\left(y_N-\sqrt{y_N^2-4g_*(N)}\right)
=
\frac{2g_*(N)}{y_N+\sqrt{y_N^2-4g_*(N)}}
\to0
\]
in probability. Therefore
\[
\widehat\sigma_i^{\mathrm{BBP}}(N)-\sigma_i^*(N)\to0
\qquad\text{in probability},
\]
which proves (iii).
\end{proof}

\subsection{Proofs of the abstract-generalization corollaries and theorems}
\label{app:proofs_abstract_generalizations}

The results below are the abstract additive-perturbation analogues of the Gaussian-instance proofs of Appendix~\ref{app:proofs_main}.  Each is obtained by substituting the abstract logit-perturbation lemma (Lemma~\ref{lem:remove_R_general}) for the Gaussian logit perturbation bound (Lemma~\ref{lem:remove_R}) and applying the abstract Assumptions~\ref{as4}--\ref{as6} in place of Gaussian Assumptions~\ref{as2}--\ref{as3}.

\begin{proof}[Proof of Theorem~\ref{theorem_reduction_of_loss_general}]
For \(s\in T\), let \(X(\alpha,s)\) denote the vector of pre-softmax logits. By Lemma~\ref{lem:remove_R_general}, for each fixed \(s\in T\),
\[
\PP\big(\|X(\alpha_S,s)-X(\alpha_W,s)\|_\infty\le a_\phi(N,s)\big)\ge 1-d_2(N),
\]
where \(a_\phi(N,s)=h_\phi(s)d_1(N)\). The single-sample cross-entropy is \(2\)-Lipschitz in \(\|\cdot\|_\infty\), so on this event \(|\ell_s(X(\alpha_S,s))-\ell_s(X(\alpha_W,s))|\le 2a_\phi(N,s)\). A union bound over \(s\in T\) gives
\[
\PP\!\left(|L_{\cancel{\mathrm{reg}}}(\alpha_S)-L_{\cancel{\mathrm{reg}}}(\alpha_W)|\le\tfrac{2}{|T|}\sum_{s\in T}a_\phi(N,s)\right)\ge 1-|T|d_2(N).
\]
Since \(a_\phi(N,s)\to0\), \(L_{\cancel{\mathrm{reg}}}(\alpha_S)-L_{\cancel{\mathrm{reg}}}(\alpha_W)\to0\) in probability. Combined with the exact regularizer identity \(\|W\|_F^2-\|S\|_F^2=\|R\|_F^2+2\langle S,R\rangle_F=\|R\|_F^2+o_{\PP}(1)\), this gives
\(
L(\alpha_S)=L(\alpha_W)-\lambda_{\mathrm{reg}}\|R\|_F^2+o_{\PP}(1).
\)
\end{proof}

\begin{proof}[Proof of Corollary~\ref{cor:scaled_reduction_of_loss_general_finiteN}]
Let \(E_N^{\mathrm{logit}}\) be the samplewise logit event \(\|X(\alpha_S,s)-X(\alpha_W,s)\|_\infty\le a_\phi(N,s)\) for all \(s\in T\). By Lemma~\ref{lem:remove_R_general} and a union bound, \(\PP(E_N^{\mathrm{logit}})\ge1-|T|d_2(N)\). On this event the samplewise perturbation argument with \((1-\vartheta)R\) in place of \(R\) gives, for every \(\vartheta\in[0,1]\),
\[
|L_{\cancel{\mathrm{reg}}}(\alpha_{W(\vartheta)})-L_{\cancel{\mathrm{reg}}}(\alpha_W)|\le (1-\vartheta)\tfrac{2}{|T|}\sum_{s\in T}a_\phi(N,s).
\]
Let \(E_N^{\mathrm{cross}}(\eta):=\{|\langle S,R\rangle_F|\le\eta\}\) and \(E_{N,\eta}:=E_N^{\mathrm{logit}}\cap E_N^{\mathrm{cross}}(\eta)\). Then \(\PP(E_{N,\eta})\ge 1-|T|d_2(N)-\PP(|\langle S,R\rangle_F|>\eta)\). On \(E_{N,\eta}\), the exact regularizer identity \(\|W\|_F^2-\|W(\vartheta)\|_F^2=(1-\vartheta^2)\|R\|_F^2+2(1-\vartheta)\langle S,R\rangle_F\) implies
\[
L(\alpha_{W(\vartheta)})\le L(\alpha_W)-\lambda_{\mathrm{reg}}(1-\vartheta^2)\|R\|_F^2+(1-\vartheta)\tfrac{2}{|T|}\sum_{s\in T}a_\phi(N,s)+2\lambda_{\mathrm{reg}}(1-\vartheta)\eta
\]
for every \(\vartheta\in[0,1]\), which is exactly the claim.
\end{proof}

\begin{proof}[Proof of Corollary~\ref{cor:general_BGN_persistent_spikes}]
Since \(\theta_c(N)\to0\) and \(\sigma_i^*(N)\to\bar\sigma_i>0\) in probability, \(\PP(\sigma_i^*(N)>\theta_c(N))\to1\) for each \(1\le i\le k\), which proves (i). On the same high-probability event, the assumed BGN-type deformation regime places the \(i\)-th spike above the critical threshold, so the corresponding singular value cluster of \(W^*(N)\) lies above the bulk edge. If the gap condition in (ii) holds for a spike, Wedin's theorem gives \(\|\widehat P_i(N)-P_i^*(N)\|_{\mathrm{op}}\le C\|R^*(N)\|_{\mathrm{op}}/\Delta_i(N)\to0\) in probability, for the left rank-one projectors (and similarly for the right projectors), proving (ii). Applying Theorem~\ref{thm:abstract_inverse_D} with \(\theta_i(N)=\sigma_i^*(N)\) then gives \(\widehat\sigma_i^{(D)}(N)-\sigma_i^*(N)\to0\) in probability for each fixed \(1\le i\le k\), which is (iii).
\end{proof}

\begin{proof}[Proof of Corollary~\ref{cor:growing_set_abstract_loss}]
By Lemma~\ref{lem:remove_R_general} and a union bound over \(s\in T_N\), with probability at least \(1-|T_N|d_2(N)\),
\[
\|X(\alpha_S,s)-X(\alpha_W,s)\|_\infty\le a_\phi(N,s)
\qquad\forall s\in T_N.
\]
On this event, the \(2\)-Lipschitz property of cross-entropy gives
\[
\left|
L_{\cancel{\mathrm{reg}},T_N}(\alpha_S)
-
L_{\cancel{\mathrm{reg}},T_N}(\alpha_W)
\right|
\le
\frac{2}{|T_N|}\sum_{s\in T_N}a_\phi(N,s).
\]
The right-hand side is \(o_{\PP}(1)\), and the event has probability tending to one because \(|T_N|d_2(N)\to0\). The regularized identity
\[
\|W\|_F^2-\|S\|_F^2
=
\|R\|_F^2+2\langle S,R\rangle_F
\]
and the cross-term hypothesis then yield
\[
L_{T_N}(\alpha_S)
=
L_{T_N}(\alpha_W)
-\lambda_{\mathrm{reg}}\|R\|_F^2
+o_{\PP}(1).
\]
\end{proof}

\begin{proof}[Proof of Corollary~\ref{cor:scaled_reduction_of_loss_general}]
Set \(W(\vartheta)=S+\vartheta R\). Repeating the argument above with perturbation \((1-\vartheta)R\) in place of \(R\), we get
\[
L_{\cancel{\mathrm{reg}}}(\alpha_{W(\vartheta)})-L_{\cancel{\mathrm{reg}}}(\alpha_W)=o_{\PP}(1).
\]
Also,
\[
\|W\|_F^2-\|W(\vartheta)\|_F^2
=
(1-\vartheta^2)\|R\|_F^2+2(1-\vartheta)\langle S,R\rangle_F
=
(1-\vartheta^2)\|R\|_F^2+o_{\PP}(1).
\]
Hence
\[
L(\alpha_{W(\vartheta)})
=
L(\alpha_W)-\lambda_{\mathrm{reg}}(1-\vartheta^2)\|R\|_F^2+o_{\PP}(1).
\]
The same argument, or the uniform finite-\(N\) pathwise bound of Corollary~\ref{cor:scaled_reduction_of_loss_general_finiteN}, gives the statement along any deterministic sequence \(\vartheta_N\in[0,1]\).
\end{proof}

\begin{proof}[Proof of Theorem~\ref{thm:multilayer_telescoping}]
Order the pruned layers as \(\mathcal P=\{\ell_1,\dots,\ell_m\}\). Let \(\alpha^{(j)}\) denote the network parameters after the first \(j\) replacements, so \(\alpha^{(0)}\) is the original network and \(\alpha^{(m)}=\alpha^{(|\mathcal P|)}\). By assumption, at stage \(j\) the hypotheses of Theorem~\ref{theorem_reduction_of_loss_general} apply to layer \(\ell_j\), hence
\[
L(\alpha^{(j)})
=
L(\alpha^{(j-1)})-\lambda_{\mathrm{reg}}\|R_{\ell_j}\|_F^2+o_{\PP}(1).
\]
Summing these identities over \(j=1,\dots,m\) and using that \(m=|\mathcal P|\) is fixed gives
\[
L(\alpha^{(m)})
=
L(\alpha^{(0)})-\lambda_{\mathrm{reg}}\sum_{j=1}^{m}\|R_{\ell_j}\|_F^2+o_{\PP}(1),
\]
which is the claimed finite-layer telescoping formula.

For the growing-layer version, use the finite-\(N\) single-layer bound on the intersection \(\bigcap_{\ell\in\mathcal P_N}E_{\ell,N}\), whose probability tends to one by the assumed summability of the failure probabilities. Summing the exact loss identities over the ordered layers gives the accumulated regularization decrease \(-\lambda_{\mathrm{reg}}\sum_{\ell\in\mathcal P_N}\|R_\ell\|_F^2\) plus the sum of the cross-entropy perturbation errors and the accumulated cross term. The two summability hypotheses make these remainders \(o_{\PP}(1)\), yielding the same formula with \(\mathcal P_N\).
\end{proof}

\begin{proof}[Proof of Theorem~\ref{main_result_reducing_noise_general}]
Let \(b_N:=(2/|T|)\sum_{s\in T}a_\phi(N,s)\to0\).
Since \(\langle S^*(N),R^*(N)\rangle_F=o_{\PP}(1)\), choose a deterministic sequence \(\eta_N\downarrow0\) such that \(\PP(|\langle S^*(N),R^*(N)\rangle_F|>\eta_N)\to0\).
By Corollary~\ref{cor:scaled_reduction_of_loss_general_finiteN} applied at the limit point, there is an event \(E_N\) with \(\PP(E_N)\to1\) such that on \(E_N\), for every \(\vartheta\in[0,1]\),
\[
L(\alpha_\vartheta^*(N))-L(\alpha^*(N))
\le
-\lambda_{\mathrm{reg}}(1-\vartheta^2)\|R^*(N)\|_F^2
+(1-\vartheta)b_N
+2\lambda_{\mathrm{reg}}(1-\vartheta)\eta_N.
\]
Write \(\vartheta=1-h\) and divide by \(h>0\):
\[
\frac{L(\alpha_{1-h}^*(N))-L(\alpha^*(N))}{h}
\le
-\lambda_{\mathrm{reg}}(2-h)\|R^*(N)\|_F^2
+b_N+2\lambda_{\mathrm{reg}}\eta_N
\quad\forall h\in(0,1].
\]
On the high-probability event where the one-sided directional stationarity hypothesis also holds, taking \(\liminf_{h\downarrow0}\) yields
\(
2\lambda_{\mathrm{reg}}\|R^*(N)\|_F^2
\le
b_N+\epsilon_N+2\lambda_{\mathrm{reg}}\eta_N,
\)
so
\[
\|R^*(N)\|_F^2
\le
(b_N+\epsilon_N)/(2\lambda_{\mathrm{reg}})+\eta_N+o_{\PP}(1).
\]
Since \(\eta_N\to0\), this proves the quantitative estimate; the vanishing-in-probability statement follows whenever \(b_N+\epsilon_N\to0\) in probability. The pathwise conclusion under the convergence/continuity hypotheses follows because \(\|R(t,N)\|_F\to\|R^*(N)\|_F\) in probability and \(\PP(\|R^*(N)\|_F=\epsilon)=0\).
\end{proof}

\bibliographystyle{plainnat}
\bibliography{rmt_vit_pruning}
\end{document}
